% !TeX program = pdflatex
% Компиляция: bash build_paper.sh  (pdflatex прогоняется до схождения
% перекрёстных ссылок; двух проходов не хватает, нужно три)
%
% ВЕРСИЯ 2. Работа переведена с моментного вывода схемы B на ТОЧНЫЙ
% условный вывод: K_b = N_b - 1 | D_b ~ Poisson(lambda_b D_b).
% Все числовые значения — вывод analysis_v2.py при зерне 20260914
% (results_v2.txt, теги [S0]...[S12]).
%
% ВЕРСИЯ 3 (редакция 24.09.2026). Изменения по существу:
%  * предложение о точном условном законе переформулировано как УСЛОВНОЕ
%    ПРАВДОПОДОБИЕ: D_b достаточна для W_b при фиксированной kappa
%    (факторизация); показано совпадение с интегрированным правдоподобием
%    при мере dW_b (мера Хаара сдвигов) и связь с задачей Неймана — Скотта;
%  * условие (A1) сформулировано для СЛУЧАЙНОГО окна, независимого от процесса;
%  * «доказуемая двусторонняя вилка» [(M-B)/S, (M-1)/S] ОТОЗВАНА: без (A1)
%    строго держится лишь верхняя граница (тождество Вальда); правило
%    остановки «после долгой засухи» даёт смещение (M-B)/S вверх
%    (новый тег [S13], скрипт stopping_rules.py);
%  * исправлен множитель смещения МНК при схеме B (две разные регрессии);
%  * неверная ссылка на Pomerance 1981 (псевдопростые) заменена на
%    Math. Intelligencer 3 (1981) 97-105;
%  * мелкие численные уточнения (покрытие Гарвуда, p однородности,
%    (1+mu)e^{-mu} при mu=16, число событий в расчёте мощности);
%  * НОВЫЙ раздел sec:second: член второго порядка самой эвристики ЛПВ,
%    lambda_b(t) = (kappa/ln b)(1 + c_b/t), точный условный вывод при заданном c
%    (lpw_second_order.py, теги [S14]-[S17], вывод в results_v3.txt);
%  * закрыты все маркеры \recompute (модель eq:logcorr, B=20, исключение по
%    одному основанию);
%  * НОВЫЙ подраздел sec:frontiers: документированные фронты b=2,10,22,26,
%    эмпирическая проверка (A1) и гибридная оценка (frontiers.py, тег [S18]);
%    отмечены PRP R_5(4939471), R_5(5154509) 2024 г. вне OEIS.
\documentclass[11pt,a4paper]{article}

% cmap ДОЛЖЕН идти до fontenc: он встраивает в PDF таблицы ToUnicode, без
% которых кириллица не извлекается из готового файла — текст нельзя ни найти
% поиском, ни скопировать, ни разобрать системам индексации (РИНЦ, Crossref).
\usepackage{cmap}
\usepackage[T2A]{fontenc}
\usepackage[utf8]{inputenc}
\usepackage[english,russian]{babel}
\usepackage{amsmath,amssymb,amsthm}
\usepackage{booktabs}
\usepackage{graphicx}
% Рисунки генерирует analysis_v2.py и кладёт в figs/.
\graphicspath{{figs/}{./}}
\usepackage[margin=25mm]{geometry}
% xurl ДО hyperref: разрешает переносить длинные URL в любом месте. Без него
% ссылки вроде t5k.org/mersenne/heuristic.html вылезают за правое поле.
\usepackage{xurl}
\usepackage[colorlinks=true,allcolors=blue]{hyperref}
% Запас на растяжение межсловных пробелов: в русском тексте с большим числом
% неразрывной формульной вёрстки без него остаются переполненные строки.
\emergencystretch=3em
\usepackage{xcolor}
\usepackage[expansion=false]{microtype}
% Подписи в принятом отечественном виде: «Рис. 1. Название», «Таблица 1. Название»
% (вместо «Рис. 1: Название»).
\usepackage{caption}
\captionsetup{labelsep=period}
% п. 4.1.5 ГОСТ Р 7.0.7-2021: подписи к иллюстративному материалу приводят на
% языке текста статьи и, как правило, повторяют на английском языке.
% Номер берётся из счётчика, уже увеличенного \caption, поэтому русская и
% английская подписи всегда несут один и тот же номер.
\newcommand{\encaptab}[1]{\\[3pt]\foreignlanguage{english}{%
  \textbf{Table~\thetable.} #1}}
\newcommand{\encapfig}[1]{\\[3pt]\foreignlanguage{english}{%
  \textbf{Figure~\thefigure.} #1}}

% Маркер величин, которые НЕЛЬЗЯ получить из троек (N_b, n_min, n_max) и
% которые требуют полных списков индексов: они пересчитываются
% analysis_v2.py при наличии b-файлов (ключ --oeis-dir). Перед подачей
% запустить скрипт с данными и заменить маркеры на числа; команду
% переопределить пустой, чтобы убедиться, что маркеров не осталось.
\newcommand{\recompute}[1]{\textcolor{red}{[пересчитать: #1]}}

\theoremstyle{plain}
\newtheorem{proposition}{Предложение}
\newtheorem{corollary}{Следствие}
\theoremstyle{remark}
\newtheorem{remark}{Замечание}

\newcommand{\Rb}{R_b}
\newcommand{\eg}{e^{\gamma}}
\newcommand{\emg}{e^{-\gamma}}
\DeclareMathOperator{\Gam}{Gamma}
\DeclareMathOperator{\Pois}{Poisson}

\title{Схема наблюдения в статистике обобщённых репьюнитных\\
простых: точный условный критерий константы\\
Ленстры --- Померанса --- Вагстаффа}
\author{Тимофей Сергеевич Докучаев\thanks{ORCID: \url{https://orcid.org/0009-0006-0510-5225}.
Контактный адрес: \texttt{timofei93@yandex.ru}}\\
\small Независимый исследователь, Иркутск, Россия}
% Дата фиксирована: \today менялась бы при каждой пересборке.
\date{27 августа 2026~г.}

\begin{document}

% Тип статьи (п. 4.4) и индекс УДК (п. 4.5) ГОСТ Р 7.0.7-2021:
% приводят в начале статьи отдельными строками слева.
\noindent Научная статья\\
\noindent УДК 511.3:519.234
\medskip

% article.cls открывает заголовок командой \newpage — причём дважды: один раз
% в самом \maketitle, второй внутри \@maketitle. Строки «Научная статья» и
% «УДК» к этому моменту уже набраны, поэтому \newpage сбрасывал их отдельной
% страницей, а заглавие уезжало на вторую. Локально обезвреживаем \newpage
% ровно на время \maketitle: строки остаются перед заглавием, как требуют
% пп. 4.4 и 4.5 ГОСТ Р 7.0.7-2021, и всё умещается на первой странице.
\let\RHnewpage\newpage
\let\newpage\relax
\maketitle
\let\newpage\RHnewpage

\begin{abstract}
Эвристика Ленстры --- Померанса --- Вагстаффа (ЛПВ) предсказывает, что
$\#\{n\le x:\Rb(n)\ \text{просто}\}\sim(\eg/\ln b)\ln x$. С 1993~года эту
универсальность проверяют через <<нормированную частоту встречаемости>>,
получаемую регрессией $\ln n_k$ на~$k$.

Во-первых, подозрение о смещении этой частоты правым цензурированием
\emph{неверно}: и концевая форма, и наклон МНК несмещены точно при схеме
<<до $N$-го события>> и с точностью до экспоненциально малого по $\mu$ члена
при фиксированном фронте (не более $0{,}73\%$ для концевой формы и $2{,}3\%$ для
МНК при $\mu\ge7$; предложение~\ref{prop:unbiased}); смещена лишь обратная
величина --- неравенство Йенсена. Дефект иной: варианты оценки расходятся до
$29{,}5\%$ --- в $1{,}7$ раза больше медианного отклонения от $\emg$, --- и
ярлык <<удачливости>> основания переворачивается от смены формы.

Во-вторых, объединение $B=15$ оснований даёт $M=218$ событий против $51$ у
одних простых Мерсенна, но употребительный пивот $\kappa S\sim\Gam(M,1)$ точен
лишь при наблюдении <<до $N$-го события>>, а данные получены при фиксированном
фронте со случайным числом находок, причём сам фронт не записан. Неизвестные
фронты --- мешающие параметры, по одному на основание, и профильная оценка
$M/S$ несостоятельна по Нейману --- Скотту. Главный результат --- точное
условное правдоподобие схемы B (предложение~\ref{prop:cond}): положение
последнего события $D_b$ достаточно для фронта, и условно на нём величина
$N_b-1$ распределена как $\Pois(\lambda_bD_b)$ \emph{точно} и \emph{независимо от
положения фронта}. Отсюда $M-B\mid D\sim\Pois(\kappa S)$: правильным пивотом
схемы~B является пуассоновский, а не гамма, правильной объединённой оценкой ---
$(M-B)/S$, а не $(M-1)/S$ (обусловливание стоит одного события \emph{на каждое
основание}, как одна степень свободы на группу в задаче Неймана --- Скотта), и
доверительные интервалы получаются в замкнутой
форме без симуляции. При номинальном $5\%$ уровне старый критерий отвергает
верную гипотезу в $11$--$17\%$ случаев в зависимости от положения фронта. Новый
вывод: $\hat\kappa=1{,}838$, точный $95\%$ интервал $[1{,}59;\,2{,}11]$,
$p=0{,}67$ против $\eg\approx1{,}781$; сообщавшееся десятипроцентное превышение
над ЛПВ есть артефакт схемы наблюдения.

В-третьих, вывод опирается на одно допущение: фронт выбирался без оглядки на
находки. При его нарушении оценки расходятся на $6{,}9\%$ и больше, и данные
этого не различают. Без допущения строго держится лишь верхняя граница (тождество
Вальда), а правило <<остановиться после долгой полосы без находок>> смещает и
$(M-B)/S$ вверх. Четыре документированных фронта впервые позволяют проверить
допущение по данным: они лежат ближе к последним находкам, чем оно предсказывает,
то есть опасного направления нет. Отсюда требование публиковать фронты $L_b$. Наконец,
условный мультиномиальный критерий однородности ($\Lambda=7{,}13$, точное
$p=0{,}935$) различает лишь разброс $\kappa_b$ свыше $\pm37\%$, а объединённая
выборка --- отклонения от $\eg$ свыше $21\%$.

В-четвёртых, сама эвристика ЛПВ предсказывает при конечных $n$ интенсивность
$(\kappa/\ln b)(1+c_b/t)$ с $c_b\in[\ln2;\,1{,}74]$, поскольку делители $\Rb(p)$ не
меньше $2p+1$. Точный условный вывод переносится на эту модель. Член $c_b$ данными не
определяется, но сдвигает оценку $\kappa$ на $8$--$16\%$ --- сильнее, чем выбор схемы
наблюдения. С ним ЛПВ согласуется с данными ($p=0{,}63$--$0{,}74$ при $c=\ln2$,
$n>10$), исчезает тренд при усечении и объясняется фольклорная <<медленная
сходимость снизу>>.

\medskip
\noindent\textbf{Ключевые слова:} репьюнит-простые, обобщённые репьюниты,
числа Мерсенна, гипотеза Ленстры --- Померанса --- Вагстаффа, постоянная
Эйлера --- Маскерони, схема наблюдения, цензурирование типа~I, обратная
выборка, пуассоновский точечный процесс, условное правдоподобие, мешающие
параметры, задача Неймана --- Скотта, достаточная
статистика, смещение оценки, мощность критерия, вычислительная теория чисел

\medskip
\noindent\textbf{MSC 2020:} Primary 11Y11; Secondary 11N05, 62F03, 62B05,
62N01, 60G55, 11N80.

\medskip
\noindent\textbf{Благодарности:} автор благодарит проект Great Internet Mersenne
Prime Search (GIMPS) за предоставление данных о простых числах Мерсенна через
открытый отчёт PrimeNet (\url{https://www.mersenne.org/}); значения индексов
Мерсенна получены оттуда 27~августа 2026~года. Автор признателен участникам
OEIS, которые много лет поддерживают последовательности обобщённых репьюнитных
простых. Автор благодарит Jaume Gotanegra (MathOverflow,
\url{https://mathoverflow.net/questions/514719}) за замечание о том, что условное
правдоподобие~\eqref{eq:condlaw} получается также интегрированием по фронту с
мерой $dW_b$ (замечание~\ref{rem:integrated}).

\medskip
\noindent\textbf{Финансирование:} исследование выполнено без внешнего
финансирования.
\end{abstract}

\begin{otherlanguage}{english}

% Англоязычный блок: заглавие, автор, аффилиация, аннотация, ключевые слова.
% Требуется журналами перечня ВАК и всеми индексируемыми в Scopus/WoS.
\begin{center}
{\large\bfseries The observation scheme in the statistics of generalized\\
repunit primes: an exact conditional test of the\\
Lenstra--Pomerance--Wagstaff constant}

\medskip
{\bfseries Timofei Sergeevich Dokuchaev}

\smallskip
% п. 4.9 ГОСТ Р 7.0.7-2021: организация, город, страна, электронный адрес
% (без слова «e-mail») и ORCID разделяют запятыми; точку в конце не ставят.
{\small Independent researcher, Irkutsk, Russia,\\
\texttt{timofei93@yandex.ru}, \url{https://orcid.org/0009-0006-0510-5225}}
\end{center}

\begin{abstract}
\noindent The Lenstra--Pomerance--Wagstaff (LPW) heuristic
predicts $\#\{n\le x: (b^n-1)/(b-1)\ \text{prime}\}\sim(e^{\gamma}/\ln b)\ln x$.
Since 1993 it has been tested via a ``normalised occurrence frequency'' from
regressing $\ln n_k$ on $k$.

First, the suspicion that right-censoring biases this frequency downwards is
\emph{false}: both the endpoint form and the OLS slope are exactly unbiased
under inverse sampling, and unbiased up to a term exponentially small in $\mu$
under a fixed frontier (at most $0.73\%$ for the endpoint form and $2.3\%$ for
OLS when $\mu\ge7$); only the reciprocal is biased --- Jensen's inequality. The
genuine defect is that competing forms of the estimator disagree by up to
$29.5\%$ on identical data, and the ``lucky base'' label flips between them.

Second, pooling $B=15$ bases yields $M=218$ events against $51$ for Mersenne
primes alone, but the customary pivot $\kappa S\sim\Gam(M,1)$ is exact only
under inverse sampling, whereas the data arise from a fixed frontier with a
random number of finds, and the frontier itself is not recorded. The unknown
frontiers are nuisance parameters, one per base, and the profile estimator $M/S$
is inconsistent in the sense of Neyman and Scott. Our main result is an exact
conditional likelihood for the latter scheme (Proposition~\ref{prop:cond}): the
position $D_b$ of the last event is sufficient for the frontier, and
conditionally on it $N_b-1$ is \emph{exactly} $\Pois(\lambda_bD_b)$, a law that
does not involve the frontier at all. Hence
$M-B\mid D\sim\Pois(\kappa S)$: the correct pivot under a fixed frontier is
Poisson rather than Gamma, the correct pooled estimator is $(M-B)/S$ rather
than $(M-1)/S$ --- conditioning costs one event \emph{per base}, as one degree
of freedom per group in the Neyman--Scott problem ---
and confidence intervals are available in closed form without simulation.
At the nominal $5\%$ level the old test rejects a true hypothesis in
$11$--$17\%$ of cases depending on the frontier. The revised conclusion is
$\hat\kappa=1.838$, exact $95\%$ interval $[1.59,\,2.11]$, $p=0.67$ against
$e^{\gamma}\approx1.781$; the reported ten-percent excess is an artefact of the
scheme.

Third, the inference rests on a single assumption: the frontier was chosen
without regard to the finds. If it fails, the estimators differ by $6.9\%$ or
more, and the data cannot tell which is right. Without the assumption only an
upper bound holds rigorously (Wald's identity), and a rule such as ``stop after
a long stretch without finds'' biases even $(M-B)/S$ upwards. Four documented
frontiers allow the first empirical check of the assumption: they lie closer to the
last finds than it predicts, so the dangerous direction is not supported. Hence our
call to publish the frontiers $L_b$. Finally, the exact
conditional multinomial homogeneity test ($\Lambda=7.13$, $p=0.935$) detects
only a spread of $\kappa_b$ above $\pm37\%$, the pooled sample only deviations
above $21\%$.

Fourth, the LPW heuristic itself predicts a finite-$n$ intensity
$(\kappa/\ln b)(1+c_b/t)$ with $c_b\in[\ln2,\,1.74]$, because the divisors of $R_b(p)$
are at least $2p+1$, and the exact conditional inference carries over to this model.
The data do not determine $c_b$, yet it moves the estimate of $\kappa$ by
$8$--$16\%$ --- more than the choice of observation scheme. With it, LPW fits the
data ($p=0.63$--$0.74$ for $c=\ln2$, $n>10$), the truncation trend disappears, and
the folklore ``slow convergence from below'' is explained.

\medskip
\noindent\textbf{Keywords:} repunit primes, generalized repunits, Mersenne
primes, Lenstra--Pomerance--Wagstaff conjecture, Euler--Mascheroni constant,
observation scheme, Type-I censoring, inverse sampling, Poisson point process,
conditional likelihood, nuisance parameters, Neyman--Scott problem, sufficient
statistic, estimator bias, statistical power, computational number theory

\medskip
\noindent\textbf{Acknowledgments:} the author thanks the Great Internet Mersenne
Prime Search (GIMPS) project for providing the Mersenne prime data through the
open PrimeNet report (\url{https://www.mersenne.org/}); the Mersenne exponents
were retrieved on 27 August 2026. The author is grateful to the OEIS contributors
who have maintained the generalized repunit prime sequences for many years, and
Jaume Gotanegra (MathOverflow, \url{https://mathoverflow.net/questions/514719})
for pointing out that
the conditional likelihood~\eqref{eq:condlaw} is also obtained by integrating
out the frontier against $dW_b$ (Remark~\ref{rem:integrated}).

\medskip
\noindent\textbf{Funding:} this research received no external funding.
\end{abstract}
\end{otherlanguage}

\section{Введение}

Для целого основания $b\ge2$ обобщённый репьюнит есть
\begin{equation}\label{eq:Rb}
\Rb(n)=\frac{b^{n}-1}{b-1},
\end{equation}
причём $\Rb(n)$ может быть простым только при простом~$n$. При $b=2$ это числа Мерсенна
$M_p=2^{p}-1$. Опираясь на эвристики Вагстаффа~\cite{wagstaff1983} и
Померанса~\cite{pomerance1981}, гипотеза ЛПВ утверждает, что
\begin{equation}\label{eq:LPW}
\pi_b(x)=\#\{n\le x:\ \Rb(n)\ \text{простое}\}\ \sim\ \frac{\eg}{\ln b}\,\ln x,
\end{equation}
где $\gamma$~--- постоянная Эйлера~--- Маскерони; подробный вывод эвристики для
$b=2$ изложен в справочнике PrimePages~\cite{t5k}. Основание входит
в~\eqref{eq:LPW} лишь через множитель $1/\ln b$: константа $\eg\approx1{,}7811$
объявляется универсальной.
Дабнер~\cite{dubner1993} предложил проверять эту универсальность эмпирически по разным
основаниям, а Бурделэ~\cite{bourdelais2009} с тех пор ведёт текущие оценки нормированной
частоты, публикуя их в комментариях к записям OEIS~\cite{oeis}. Родственные вопросы для
одной лишь последовательности Мерсенна недавно обсуждались в~\cite{wagstaff2025}.
Близка по духу и работа~\cite{dominguez2026}, где у показателей Мерсенна ищут
вторичную структуру, связанную с разложением $p-1$; методически она устроена так же,
как наш~\S\ref{sec:struct} --- стратифицированное условное правдоподобие плюс
проверка значимости, --- но проверяет структуру \emph{внутри} одного основания,
тогда как мы проверяем её \emph{между} основаниями.

Вокруг этих оценок сложился интерпретационный фольклор. Основание~13 описывается как
<<удачливое>>, основание~7 --- как <<самое неудачливое>>; отмечается, что нормированные
частоты лежат ниже предполагаемого предела $\emg\approx0{,}5615$ и медленно сходятся
к нему снизу. Настоящая работа разбирает эти утверждения, но приходит к выводам,
отличным от тех, которых естественно ожидать.

\subsection{Вклад работы}

\begin{enumerate}
\item\textbf{Употребительная оценка частоты не смещена.} Схема, при которой интервал
подгонки заканчивается наибольшим известным индексом, действительно обусловливает
оценку тем, что последнее наблюдение является событием, и достаточная статистика
распределена как $\Gam(N,\lambda_b)$ (предложение~\ref{prop:gamma}). Но отсюда
\emph{не} следует смещение нормированной частоты: и концевая форма
$D/(N\ln b)$, и наклон МНК несмещены точно при схеме~A и с точностью до
множителя $1-(1+\mu)e^{-\mu}$ при схеме~B (предложение~\ref{prop:unbiased},
таблица~\ref{tab:mc}); при наблюдаемых $\mu\approx N_b\ge7$ дефицит не
превосходит $0{,}73\%$ для концевой формы и $2{,}3\%$ для МНК без начала отсчёта.
Смещена только обратная величина
$\hat\kappa=1/\hat G$, в $N/(N-1)$ раз, и это неравенство Йенсена, а не цензурирование.
Несмещённые оценки $\kappa$ и $1/\kappa$ не могут быть взаимно обратными; всякая
<<поправка>>, чинящая одну, портит другую (\S\ref{sec:notbias}).

\item\textbf{Настоящий дефект --- неустойчивость к выбору формы оценки.}
В обращении несколько численно различных вариантов. На одних и тех же данных они
расходятся до $29{,}5\%$ (основание~18) --- в $1{,}7$ раза больше медианного
отклонения оценок от $\emg$ ($17{,}6\%$), которое и толкуют содержательно; у двух
оснований из пятнадцати выбор формы весит больше, чем всё их отклонение от
предсказания ЛПВ. Основание~18 при концевой форме --- третье
по <<неудачливости>> из пятнадцати, а при МНК --- десятое, ниже медианы. Ярлык
описывает выбор оценки, а не свойства простых (\S\ref{sec:instability}).

\item\textbf{Главный результат: точное условное правдоподобие схемы B.} Пивот
$\kappa S\sim\Gam(M,1)$ точен при схеме <<наблюдаем до $N$-го события>>, где $N$
задано планом. Данные же получены при фиксированном фронте: поиск покрыл
$(t_0,\ln L_b]$, а число находок случайно. Незаписанные ширины окон $W_b$ ---
мешающие параметры, число которых растёт вместе с числом оснований; их
совместная ОМП $\hat W_b=D_b$ приводит к несостоятельной профильной оценке
$M/S$ --- в точности ситуация Неймана --- Скотта~\cite{neymanscott1948}. Мы
показываем (предложение~\ref{prop:cond}), что $D_b$ достаточна для $W_b$ при
фиксированной $\kappa$ и что условно на $D_b$ величина $N_b-1$ есть
$\Pois(\lambda_bD_b)$, причём \emph{ширина окна $W_b$ в это распределение не
входит}. Отсюда $M-B\mid D\sim\Pois(\kappa S)$ --- точный пуассоновский пивот,
а $(M-B)/S$ --- оценка условного максимального правдоподобия в смысле
Андерсена~\cite{andersen1970}, точно условно несмещённая; она же получается
интегрированием по $W_b$ с мерой $dW_b$ (замечание~\ref{rem:integrated}). Ни моментных
приближений, ни экспоненциально малых остатков, ни симуляционной инверсии
интервалов не требуется (\S\ref{sec:scheme}). Само различие схем классическое, а
несмещённость $(N-1)/D$ при схеме~A --- учебниковый факт, и мы на них не
претендуем; новизна в том, что фронт \emph{не документирован}, поэтому экспозицию
приходится восстанавливать из самих находок, и поправка оказывается
пропорциональна числу оснований (замечание~\ref{rem:priority}).

\item\textbf{Численное следствие.} Точный условный вывод даёт
$\hat\kappa=1{,}838$ при $95\%$ интервале $[1{,}594;\,2{,}109]$ и $p=0{,}672$:
превышение над ЛПВ исчезает. Старый гамма-пивот при схеме~B отвергает верную
гипотезу в $11$--$17\%$ случаев (\S\ref{sec:calib}); разброс зависит от того, как
далеко ушёл фронт, и именно поэтому приводится диапазоном.

\item\textbf{Единственное допущение и его цена.} Весь вывод держится на
условии (A1): фронт выбирался без оглядки на находки. Оценки $(M-1)/S$ и
$(M-B)/S$ построены на одних и тех же статистиках $(N_b,D_b)$ и различаются на
$6{,}9\%$; какая из них верна, определяется не данными, а тем, как они
собирались. Без (A1) строго держится лишь \emph{верхняя} граница: по тождеству
Вальда $\mathbb E[M]\ge\kappa\,\mathbb E[S]$ при любом правиле остановки.
Нижней границы нет: правило <<прекратить поиск после долгой полосы без находок>>
смещает вверх и $(M-B)/S$ (замечание~\ref{rem:bracket}, таблица~\ref{tab:stop}).
Попытка обойти выбор через заведомо несмещённые по-основанию куски проваливается
(\S\ref{sec:noescape}). Четыре документированных фронта ($b=2,10,22,26$) дают
первую эмпирическую проверку (A1). Фронты оказываются ближе к последним находкам,
чем предсказывает (A1) ($p=0{,}10$), то есть отклонение, если оно есть, направлено в
безопасную сторону $m<1$ (\S\ref{sec:frontiers}). Единственный радикальный выход ---
публиковать фронты $L_b$.

\item\textbf{Удачливые основания: чего критерий не может.} Обусловливание на
$K=M-B$ делает нулевое распределение критерия однородности точным
(мультиномиальным), без апелляции к асимптотике $\chi^2$. Получаем
$\Lambda=7{,}13$ при $14$ степенях свободы, точное $p=0{,}935$. Но мощность
такова, что критерий различает лишь разброс $\kappa_b$ порядка $\pm37\%$
(\S\ref{sec:homog}). Корректный вывод --- <<базозависимые частоты данными не
требуются>>, а не <<их не существует>>.

\item\textbf{Член второго порядка самой эвристики.} Эвристика ЛПВ предсказывает
при конечных $n$ не постоянную интенсивность, а
$\lambda_b(t)=(\kappa/\ln b)(1+c_b/t)$ с $c_b\in[\ln2;\,1{,}74]$: делители
$\Rb(p)$ не меньше $2p+1$. Точный условный вывод переносится на эту модель в
замкнутой форме (\S\ref{sec:second}). Данные $c_b$ не определяют, но оценка $\kappa$
зависит от него на $8$--$16\%$, то есть сильнее, чем от схемы наблюдения. При
$c=\ln2$ и $n>10$ ЛПВ со своим членом второго порядка согласуется с данными
($p=0{,}63$--$0{,}74$), а тренд оценки при усечении исчезает. Тот же член объясняет
фольклорную <<медленную сходимость снизу>>.

\item\textbf{Явный расчёт мощности.} При двустороннем критерии объединённая
выборка различает отклонения от $\eg$ не менее $21{,}3\%$; одна
последовательность Мерсенна --- не менее $45{,}5\%$ (\S\ref{sec:power}).
Конвенция критерия указана явно, поскольку односторонний вариант даёт
заметно другой порог ($18{,}7\%$).
\end{enumerate}

В~\S\ref{sec:sameerror} разобран более частный, но поучительный сюжет: наивная регрессия
счётной функции Мерсенна порождает на вид значимый отрицательный свободный член,
который оказывается артефактом обусловливания. Мы включаем его как разобранный пример
той же ошибки в ином обличье.

\section{Модель и данные}\label{sec:data}

\subsection{Пуассоновская модель}
Положим $t=\ln n$. Эвристика~\eqref{eq:LPW} равносильна утверждению, что индексы
обобщённых репьюнитных простых по основанию~$b$ образуют однородный пуассоновский
точечный процесс на оси~$t$ с постоянной интенсивностью
\begin{equation}\label{eq:lam}
\lambda_b=\frac{\kappa}{\ln b},\qquad \text{ЛПВ:}\ \ \kappa=\eg=1{,}781072\ldots
\end{equation}
Универсальность --- это утверждение, что единственный параметр $\kappa$ управляет всеми
основаниями. Пуассоновская идеализация здесь --- рабочее допущение, а не теорема:
даже для обычных простых пуассоновский закон в коротких промежутках доказан лишь
условно, при гипотезах Харди --- Литтлвуда, и известно, что он разрушается при росте
масштаба~\cite{jha2026}. Мы принимаем~\eqref{eq:lam} как модель и всюду рассматриваем
$\kappa$ как оцениваемую величину; это сводит вопрос <<верна ли ЛПВ?>> к однопараметрическому и
делает основания объединяемыми.

\subsection{Допустимые основания}\label{sec:bases}
Основания, являющиеся точными степенями, исключаются: при $b=c^{m}$, $m>1$, число
$\Rb(n)$ допускает алгебраическую факторизацию, и процесс не имеет предполагаемого вида.
Поэтому основания $4,8,9,16$ опущены. Внутри объявленного диапазона отбор
не зависит от результатов: берутся \emph{все} основания $2\le b\le20$, не являющиеся
точными степенями, что даёт $B=15$; ни одно основание не исключено и не добавлено
по величине его оценки. Каждое из них имеет не менее восьми известных индексов
(минимум достигается при $b=18$: восемь членов A133857).

Подчеркнём, однако, что сама верхняя граница $b\le20$ --- условная круглая величина,
а не следствие критерия <<не менее восьми индексов>>: этот критерий выполняется и
дальше (A127996 для $b=21$ содержит восемь членов, A127997 для $b=22$~--- девять) и
впервые нарушается лишь при $b=23$ (A204940, шесть членов). Чтобы граница не выглядела
подобранной, мы проверили её влияние прямо. Расширение набора до всех оснований
$b\le26$, не являющихся точными степенями (добавляются $b=21,22,23,24,26$; $b=25=5^2$
исключается алгебраически), даёт $B=20$, $M=253$ и
\[
  \frac{M-1}{S}=1{,}9776,\qquad \frac{M-B}{S}=1{,}8285,
\]
то есть сдвигает основную оценку~\eqref{eq:resB} на $0{,}5\%$ --- более чем на
порядок меньше цены выбора схемы наблюдения (\S\ref{sec:scheme}). Критерий однородности
на расширенном наборе однородность не отвергает
($\Lambda=10{,}24$ при $19$ степенях свободы, точное $p=0{,}953$; тег~[S17]).
Ни один вывод работы от границы $b\le20$ не зависит; мы сохраняем её как заранее
объявленную.

\begin{remark}[Отбор по числу находок]\label{rem:selection}
Формально граница $b\le20$ объявлена заранее, но выбрана она в мире, где числа
известных членов уже опубликованы, и это стоит проверить не риторически, а по
знаку сдвига. Условный вывод \S\ref{sec:scheme} обусловливается на $N_b\ge1$ ---
это безобидно, поскольку находки есть у всех оснований; обусловливание же на
$N_b\ge7$ безобидным не было бы и смещало бы $\hat\kappa$ \emph{вверх}. Снятие
границы до $b\le26$ добавляет, в частности, основание $b=23$ с шестью членами и
сдвигает оценку \emph{вниз}, с $1{,}8381$ до $1{,}8285$. Знак сдвига тот, который
предсказывает отбор по числу находок, а величина --- на порядок меньше цены
схемы наблюдения. Мы считаем это достаточным свидетельством того, что граница
вывода не определяет.
\end{remark}

\subsection{Данные и их статус}
В таблице~\ref{tab:data} перечислены использованные последовательности,
полученные 27~августа 2026~года (сверка --- приложение~\ref{app:data}).
Четыре оговорки существенны.

\paragraph{Вероятно простые числа.} Для больших индексов большинство элементов ---
вероятно простые, установленные тестом Ферма или сильным тестом Ферма после пробного
деления. Счёт смещается вверх лишь в случае, если псевдопростое Ферма было принято за
простое, что для чисел такого размера практически не встречается. Чувствительность
проверена в~\S\ref{sec:robust}.

\paragraph{Фронты проверки.} Для $b=2$ проект GIMPS~\cite{gimps} публикует два фронта:
все показатели ниже $X_{\mathrm{ver}}=81\,648\,221$ проверены дважды, что даёт ровно $50$
простых Мерсенна, тогда как все показатели ниже $X_{\mathrm{ft}}=141\,308\,443$ проверены
хотя бы однократно, что даёт $52$. Два наибольших показателя предварительны: между ними
остаются непроверенные кандидаты. Анализ~\S\ref{sec:sameerror}, где фронт используется
явно, ограничен $50$ проверенными показателями. Объединённый анализ к этому выбору
нечувствителен (\S\ref{sec:robust}).

Оба фронта --- движущиеся величины, поэтому фиксируем момент считывания: отчёт
PrimeNet Milestones, 27~августа 2026~г. Ни один вывод от точного значения
$X_{\mathrm{ver}}$ не зависит: достаточно того, что оно лежит между
$M_{50}=77\,232\,917$ и $M_{51}=82\,589\,933$, а это верно для всех значений отчёта
за 2026~год.

\paragraph{Статус записи A000043.} Раздел DATA этой записи содержит $50$ членов и
обрывается на $77\,232\,917$: порядок следующих показателей OEIS считает недоказанным
до завершения двойной проверки. Показатели $82\,589\,933$ и $136\,279\,841$ известны и
документированы GIMPS (и упомянуты в комментариях к A000043), но формально не входят в
DATA. Мы включаем их, поскольку статистика~\eqref{eq:pivotB} использует лишь положение
последнего события, а не его порядковый номер; в~\S\ref{sec:robust} проверено, что
ограничение основания~2 пятьюдесятью дважды проверенными показателями вывода не меняет.

\paragraph{Границы поиска и полнота покрытия при $b>2$.} Для прочих оснований
датированный и атрибутированный фронт выше последней находки опубликован лишь при
$b=10$: <<Search limit is 10800000, currently>> (С.~Баталов, комментарий к A004023,
1~июля 2021~г.). В расширенном наборе к нему добавляются $b=22$ и $b=26$ с
границей $10^5$ (Р.~Прайс, комментарии к A127997 и A127999, 2012~г.). Для
остальных оснований записи OEIS содержат лишь даты находок, а участники
распределённого поиска прямо пишут, что пределы поиска по большинству оснований
неизвестны. Именно эта лакуна --- предмет~\S\ref{sec:scheme} и, как мы полагаем,
главное препятствие к количественной проверке ЛПВ; четыре документированных фронта
использованы в~\S\ref{sec:frontiers}. Отметим
отдельно: вся конструкция предполагает, что поиск покрыл $(t_0,\ln L_b]$
\emph{сплошь}. Если в просеянных диапазонах есть пропуски --- при распределённом
поиске это обычное дело, --- то $M$ занижено относительно экспозиции, и $\hat\kappa$
смещена вниз, то есть в ту же сторону, что и наш основной вывод. Проверить полноту
покрытия по опубликованным данным в общем случае нельзя. Что пропуски реальны,
показывает основание~$5$. В 2024~г. найдены вероятно простые $R_5(4\,939\,471)$ и
$R_5(5\,154\,509)$ (R.~Propper); в OEIS их нет, а сводные таблицы отмечают
промежуток $3\,300\,593$--$4\,939\,471$ как не пройденный сплошь. Для условного
вывода это существенно: закон~\eqref{eq:condlaw} требует сплошного покрытия
$(t_0,t_0+D_b)$. Поэтому мы сохраняем $n_{\max}=3\,300\,593$. Включение обоих
чисел, допустимое лишь при сплошном покрытии, дало бы $\hat\kappa=1{,}852$ вместо
$1{,}838$ (тег~[S18]). Это ещё один довод в пользу публикации не только фронтов, но
и протоколов просеивания.

\begin{table}[t]
\centering
\caption{Использованные последовательности. $N_b$ --- число событий строго правее
$t_0=\ln n_{\min}$, то есть на единицу меньше числа известных простых индексов:
наименьший индекс задаёт начало отсчёта, а не событие; $K_b=N_b-1$ --- условное
число событий, входящее в~\eqref{eq:pivotB}. Столбец $D_b$ приведён округлённым,
но все вычисления используют точные $D_b=\ln n_{\max}-\ln n_{\min}$. В итоговой
строке $M=\sum_b N_b$, $K=M-B$ и $S=\sum_b D_b/\ln b$; величина $S$ не есть сумма
столбца $D_b$. Источник --- записи OEIS, получено 27.08.2026.
\encaptab{Sequences used. $N_b$ is the number of events strictly to the right of
$t_0=\ln n_{\min}$, that is, one less than the number of known prime indices: the
smallest index sets the origin, not an event; $K_b=N_b-1$ is the conditional event
count entering~\eqref{eq:pivotB}. The $D_b$ column is rounded, but all computations
use exact $D_b=\ln n_{\max}-\ln n_{\min}$. In the last row $M=\sum_b N_b$, $K=M-B$
and $S=\sum_b D_b/\ln b$; the quantity $S$ is not the sum of the $D_b$ column.
Source: OEIS entries, retrieved 27.08.2026.}}
\label{tab:data}
\small
\begin{tabular}{rlrrrrrrr}
\toprule
$b$ & OEIS & простых & $N_b$ & $K_b$ & $n_{\min}$ & $n_{\max}$ & $D_b$ & $D_b/\ln b$\\
\midrule
2 & A000043 & 52 & 51 & 50 & 2 & 136\,279\,841 & 18{,}0371 & 26{,}0220\\
3 & A028491 & 23 & 22 & 21 & 3 & 8\,530\,117 & 14{,}8605 & 13{,}5266\\
5 & A004061 & 19 & 18 & 17 & 3 & 3\,300\,593 & 13{,}9110 & 8{,}6434\\
6 & A004062 & 17 & 16 & 15 & 2 & 3\,360\,347 & 14{,}3344 & 8{,}0002\\
7 & A004063 & 10 & 9 & 8 & 5 & 1\,264\,699 & 12{,}4409 & 6{,}3934\\
10 & A004023 & 11 & 10 & 9 & 2 & 8\,177\,207 & 15{,}2237 & 6{,}6116\\
11 & A005808 & 13 & 12 & 11 & 17 & 1\,868\,983 & 11{,}6077 & 4{,}8408\\
12 & A004064 & 14 & 13 & 12 & 2 & 769\,543 & 12{,}8604 & 5{,}1754\\
13 & A016054 & 13 & 12 & 11 & 5 & 1\,503\,503 & 12{,}6139 & 4{,}9178\\
14 & A006032 & 11 & 10 & 9 & 3 & 1\,724\,417 & 13{,}2618 & 5{,}0252\\
15 & A006033 & 10 & 9 & 8 & 3 & 639\,833 & 12{,}2704 & 4{,}5311\\
17 & A006034 & 12 & 11 & 10 & 3 & 1\,990\,523 & 13{,}4053 & 4{,}7315\\
18 & A133857 & 8 & 7 & 6 & 2 & 1\,270\,141 & 13{,}3615 & 4{,}6228\\
19 & A006035 & 11 & 10 & 9 & 19 & 209\,359 & 9{,}3074 & 3{,}1610\\
20 & A127995 & 9 & 8 & 7 & 3 & 984\,349 & 12{,}7011 & 4{,}2397\\
\midrule
\multicolumn{3}{l}{всего} & $M=218$ & $K=203$ & & & & $S=110{,}4423$\\
\bottomrule
\end{tabular}
\end{table}

\section{Что смещено, а что нет}\label{sec:notbias}

\subsection{Употребительные оценки}
Дабнер~\cite{dubner1993} и Бурделэ~\cite{bourdelais2009} характеризуют основание~$b$
нормированной частотой встречаемости, получаемой подгонкой прямой к зависимости
$\ln n_k$ от~$k$ с последующим делением на $\ln b$. Согласно~\eqref{eq:lam} ожидаемый
наклон равен $1/\lambda_b=\ln b/\kappa$, так что нормированная частота оценивает
\begin{equation}\label{eq:G}
G_b=\frac{1}{\kappa},\qquad\text{ЛПВ:}\ \ G_b=\emg=0{,}5614595\ldots
\end{equation}
В обращении несколько вариантов --- обычный МНК со свободным членом, регрессия через
начало координат и концевая форма $\hat G_b=D_b/(N_b\ln b)$, где $D_b=t_{N_b}-t_0$.
Все они обладают одной структурной особенностью: интервал подгонки заканчивается на наибольшем
известном индексе.

\subsection{Точное распределение}

\begin{proposition}[Условный закон правого конца]\label{prop:gamma}
Пусть события образуют однородный пуассоновский процесс интенсивности $\lambda$ на
$(t_0,\infty)$, и пусть $t_1<\dots<t_N$ --- его первые $N$ событий после $t_0$.
Тогда $D:=t_N-t_0$ имеет распределение $\Gam(N,\lambda)$ с плотностью
$\lambda^{N}d^{N-1}e^{-\lambda d}/(N-1)!$. При условии $D=d$ предшествующие $N-1$ событий
распределены как порядковые статистики $N-1$ независимых равномерных величин на
$(t_0,t_0+d)$.
\end{proposition}

\begin{proof}
Времена ожидания $t_1-t_0,\dots,t_N-t_{N-1}$ независимы и распределены по закону
$\mathrm{Exp}(\lambda)$, а сумма $N$ таких величин распределена как $\Gam(N,\lambda)$.
Условная равномерность есть свойство порядковых статистик пуассоновского процесса,
применённое к $(t_0,t_0+d)$.
\end{proof}

\begin{corollary}[Смещение обратной величины]\label{cor:kappa}
Оценка максимального правдоподобия для $\lambda$ по одному лишь $D$ равна
$\hat\lambda=N/D$, причём
\begin{equation}\label{eq:bias}
\mathbb E[\hat\lambda]=N\,\mathbb E[D^{-1}]=\frac{N}{N-1}\,\lambda,
\end{equation}
так как $\mathbb E[D^{-1}]=\lambda/(N-1)$ для $D\sim\Gam(N,\lambda)$ при $N\ge2$.
Оценка $\tilde\lambda=(N-1)/D$ в точности несмещена.
\end{corollary}

\subsection{Чего из следствия \ref{cor:kappa} не следует}

Естественно заключить, что раз $\hat\lambda$ завышена в $N/(N-1)$ раз, то частота
$\hat G_b$ занижена во столько же. Это неверно ни при одной из двух схем наблюдения,
разобранных ниже.

\begin{proposition}[Несмещённость частоты]\label{prop:unbiased}
\emph{(a) Схема~A.} В условиях предложения~\ref{prop:gamma}:
\begin{enumerate}
\item[\rm(i)] концевая оценка несмещена точно:
$\mathbb E\!\left[\dfrac{D}{N\ln b}\right]=\dfrac{1}{\lambda\ln b}=G_b$;
\item[\rm(ii)] наклон МНК регрессии $t_k$ на $k$ несмещён точно:
$\mathbb E[\hat\beta]=1/\lambda$;
\item[\rm(iii)] следовательно, «исправленная» величина
$\tilde G_b=\frac{N}{N-1}\hat G_b$ смещена \emph{вверх} ровно в $N/(N-1)$ раз.
\end{enumerate}
\emph{(b) Схема~B.} Пусть процесс наблюдается на окне ширины $W$, $\mu=\lambda W$;
оценки полагаются равными нулю, когда они не определены. Тогда
\begin{equation}\label{eq:BunbiasedB}
\mathbb E\!\left[\frac{D}{N}\,\mathbf 1_{\{N\ge1\}}\right]
=\mathbb E\bigl[\hat\beta_0\,\mathbf 1_{\{N\ge1\}}\bigr]
=\frac{1}{\lambda}\Bigl(1-(1+\mu)e^{-\mu}\Bigr),
\end{equation}
где $\hat\beta_0$ --- наклон МНК по точкам $(k,t_k)$, $k=0,\dots,N$, \emph{включая}
начало отсчёта. Для наклона $\hat\beta_1$ по точкам $k=1,\dots,N$ \emph{без}
начала отсчёта, определённого при $N\ge2$,
\begin{equation}\label{eq:BunbiasedOLS}
\mathbb E\bigl[\hat\beta_1\bigm|N\ge2\bigr]
=\frac{1}{\lambda}\cdot
\frac{1-\bigl(1+\mu+\mu^2/2\bigr)e^{-\mu}}{1-(1+\mu)e^{-\mu}} .
\end{equation}
Обе формы, таким образом, несмещены с точностью до экспоненциально малого по
$\mu$ члена; при $\mu=7$ дефицит составляет $0{,}73\%$ в~\eqref{eq:BunbiasedB} и
$2{,}25\%$ в~\eqref{eq:BunbiasedOLS}, при $\mu=10$ --- $0{,}05\%$ и $0{,}23\%$.
\end{proposition}

\begin{proof}
(a)(i) $\mathbb E[D]=N/\lambda$ для $D\sim\Gam(N,\lambda)$, откуда результат немедленно.
(ii) Наклон МНК есть линейная статистика $\hat\beta=\sum_k w_k t_k$ с весами
$w_k=(k-\bar k)/\sum_j(j-\bar k)^2$, для которых $\sum_k w_k=0$ и $\sum_k w_k k=1$.
Поскольку $\mathbb E[t_k]=t_0+k/\lambda$, получаем
$\mathbb E[\hat\beta]=\sum_k w_k\,(t_0+k/\lambda)=1/\lambda$.
(iii) непосредственно из (i).

(b) При схеме~B условно на $N=n\ge1$ положения событий суть порядковые статистики
$n$ независимых равномерных величин на $(t_0,t_0+W]$, поэтому $D/W\sim\mathrm{Beta}(n,1)$
и $\mathbb E[D\mid N=n]=Wn/(n+1)$, откуда $\mathbb E[D/N\mid N=n]=W/(n+1)$. Усредняя по
$N\sim\mathrm{Pois}(\mu)$ и учитывая, что $D$ определено лишь при $N\ge1$,
\[
\mathbb E\Bigl[\frac{D}{N}\Bigr]
=W\sum_{n\ge1}\frac{\mu^{n}e^{-\mu}}{n!\,(n+1)}
=\frac{W}{\mu}\sum_{n\ge1}\frac{\mu^{n+1}e^{-\mu}}{(n+1)!}
=\frac{W}{\mu}\Bigl(1-e^{-\mu}-\mu e^{-\mu}\Bigr),
\]
что и есть~\eqref{eq:BunbiasedB} для концевой формы. Для наклона МНК заметим, что
$\mathbb E[t_k\mid N=n]=t_0+kW/(n+1)$ при $k=0,\dots,n$; поскольку наклон ---
линейная статистика с весами $\sum_kw_k=0$, $\sum_kw_kk=1$, в обоих вариантах
$\mathbb E[\hat\beta\mid N=n]=W/(n+1)$. Для $\hat\beta_0$ суммирование идёт по
$n\ge1$ и даёт то же, что и выше. Для $\hat\beta_1$ оно идёт по $n\ge2$:
$\sum_{n\ge2}\frac{W}{n+1}\frac{\mu^ne^{-\mu}}{n!}
=\frac{W}{\mu}\bigl(1-(1+\mu+\mu^2/2)e^{-\mu}\bigr)$, и деление на
$\Pr(N\ge2)=1-(1+\mu)e^{-\mu}$ даёт~\eqref{eq:BunbiasedOLS}.
\end{proof}

Содержательный смысл прост: несмещённые оценки $\kappa$ и $1/\kappa$ не могут быть
взаимно обратными, поскольку $x\mapsto1/x$ строго выпукло. Несмещённая пара --- это
$(\hat G_b,\ \tilde\kappa_b)$, где $\tilde\kappa_b=(N_b-1)\ln b/D_b$; но
$\tilde\kappa_b\ne1/\hat G_b$. Всякая единая <<поправка на цензурирование>>,
исправляющая одну величину, портит другую. Таблица~\ref{tab:mc} подтверждает
предложение~\ref{prop:unbiased} симуляцией при обеих схемах.

\begin{remark}\label{rem:schemeindep}
Существенно, что пункт (b) снимает возможное возражение к первому выводу работы.
Ниже (\S\ref{sec:scheme}) мы принимаем схему~B как соответствующую способу
получения данных; если бы пункт~(a) был единственным, то весь \S\ref{sec:notbias}
опирался бы на схему, которую сама работа отвергает. Пункт~(b) показывает, что
вывод от выбора схемы не зависит: дефицит $(1+\mu)e^{-\mu}$ при $\mu\approx N_b$
составляет $0{,}73\%$ при $N_b=7$, $0{,}05\%$ при $N_b=10$ и около $2\cdot10^{-6}$ при
$N_b=16$ (для МНК без начала отсчёта~\eqref{eq:BunbiasedOLS} --- $2{,}25\%$,
$0{,}23\%$ и $1{,}4\cdot10^{-5}$), тогда как обсуждаемые в литературе отклонения ---
десятки процентов.
\end{remark}

\begin{table}[t]
\centering
\caption{Проверка предложения~\ref{prop:unbiased} методом Монте-Карло,
$60\,000$ реализаций на строку (зерно фиксировано, см.\ \texttt{analysis\_v2.py},
тег~[S1]). Целевое значение всех столбцов, кроме второго, равно $1{,}0000$.
В правой части $\mu=\lambda W$; столбец <<теория>> даёт множитель
$1-(1+\mu)e^{-\mu}$ из~\eqref{eq:BunbiasedB} для $\hat G$. Столбец
$\mathbb E[\hat\beta]\lambda$ схемы~B вычислен для наклона $\hat\beta_1$ без начала
отсчёта при условии $N\ge2$, и его теоретическое значение
даёт~\eqref{eq:BunbiasedOLS}: $0{,}9775$, $0{,}9977$ и $1{,}0000$ при $\mu=7,10,16$.
\encaptab{Monte Carlo check of Proposition~\ref{prop:unbiased}, $60\,000$
replications per row (seed fixed, see \texttt{analysis\_v2.py}, tag~[S1]). The
target value of every column except the second is $1.0000$. In the right-hand
part $\mu=\lambda W$; the column ``theory'' gives the factor
$1-(1+\mu)e^{-\mu}$ of~\eqref{eq:BunbiasedB} for $\hat G$. The scheme~B column
$\mathbb E[\hat\beta]\lambda$ is for the slope $\hat\beta_1$ without the origin,
conditional on $N\ge2$; its theoretical value is given
by~\eqref{eq:BunbiasedOLS}: $0.9775$, $0.9977$ and $1.0000$ at $\mu=7,10,16$.}}
\label{tab:mc}
\small
\begin{tabular}{rrrrrr@{\qquad}rrr}
\toprule
& \multicolumn{5}{c}{схема A ($N$ задано)} & \multicolumn{3}{c}{схема B ($\mu=\lambda W$)}\\
\cmidrule(r){2-6}\cmidrule(l){7-9}
$N=\mu$ & $\mathbb E[\hat\lambda]/\lambda$ & $N/(N-1)$
& $\mathbb E[\tilde\lambda]/\lambda$ & $\mathbb E[\hat G]\lambda$
& $\mathbb E[\hat\beta]\lambda$
& $\mathbb E[\hat G]\lambda$ & теория & $\mathbb E[\hat\beta]\lambda$\\
\midrule
7 & 1{,}1668 & 1{,}1667 & 1{,}0001 & 1{,}0003 & 1{,}0000 & 0{,}9952 & 0{,}9927 & 0{,}9783\\
10 & 1{,}1123 & 1{,}1111 & 1{,}0011 & 0{,}9997 & 1{,}0004 & 1{,}0002 & 0{,}9995 & 0{,}9971\\
16 & 1{,}0655 & 1{,}0667 & 0{,}9989 & 1{,}0012 & 1{,}0015 & 1{,}0000 & 1{,}0000 & 1{,}0006\\
22 & 1{,}0468 & 1{,}0476 & 0{,}9992 & 1{,}0005 & 0{,}9997 & 1{,}0002 & 1{,}0000 & 0{,}9995\\
51 & 1{,}0199 & 1{,}0200 & 0{,}9999 & 1{,}0001 & 1{,}0003 & 0{,}9999 & 1{,}0000 & 0{,}9999\\
\bottomrule
\end{tabular}
\end{table}

\begin{remark}
Отрицательный результат этого параграфа стоит подчеркнуть, поскольку он противоречит
интуиции: систематическое отклонение опубликованных частот вниз от $\emg$
\emph{не} объясняется схемой подгонки. Его источник --- не арифметика оценки, а
член второго порядка самой эвристики, $c_b\ln\ln x$ в счётной функции, который
занижает частоту при конечных $n$ и исчезает лишь как $1/\ln n$
(\S\ref{sec:second}); вклад в разрыв вносит и схема наблюдения
(\S\ref{sec:scheme}).
\end{remark}

\subsection{Неустойчивость к выбору формы оценки}\label{sec:instability}

Несмещённость не означает согласованности. В таблице~\ref{tab:bases} приведены
концевая форма $\hat G_b$ и оценка по МНК $\hat G_b^{\mathrm{OLS}}$ на одних и тех же
данных. Расхождения достигают $29{,}5\%$ ($b=18$), $15{,}9\%$ ($b=13$) и $10{,}3\%$
($b=11$).

Сопоставим их с тем, что в литературе толкуют содержательно, --- с отклонениями
самих оценок от $\emg$. Последние по пятнадцати основаниям имеют медиану $17{,}6\%$
(от $5{,}6\%$ при $b=20$ до $43{,}7\%$ при $b=19$). Наибольшее расхождение форм,
$29{,}5\%$, превышает эту медиану в $1{,}7$ раза. В среднем выбор формы, разумеется,
дешевле: медиана расхождений равна $4{,}1\%$, то есть около четверти медианного
отклонения. Но у двух оснований из пятнадцати ($b=18$ и $b=20$) выбор формы весит
больше, чем всё их отклонение от предсказания ЛПВ, --- и одно из них, основание~18,
фольклор как раз и обсуждает.

Следствие для фольклора прямое. Основание~18 при концевой форме имеет
$\hat G_{18}=0{,}6604$ --- третье по величине из пятнадцати (после $b=7$ с $0{,}7104$
и $b=10$ с $0{,}6612$), то есть уверенно входит в тройку <<самых неудачливых>>; при МНК
$\hat G_{18}^{\mathrm{OLS}}=0{,}4656$ --- десятое место из пятнадцати, ниже медианы
$0{,}4904$, то есть оно уже <<удачливое>>. Ярлык переворачивается без единого нового
простого числа, от одной лишь смены формы оценки: основание перемещается с третьего
места на десятое.

К этому добавляется вторая, техническая трудность МНК: его стандартные ошибки в этой задаче
неверны. Остатки суть отклонения порядковых статистик и сильно автокоррелированы;
для $b=2$ мы получаем $\rho_1=0{,}825$, а ковариационная оценка Ньюи --- Уэста с восемью
лагами увеличивает стандартную ошибку свободного члена с $0{,}447$ до $0{,}732$
(\S\ref{sec:sameerror}). Опубликованные точечные значения не сопровождаются интервалами,
и это, а не смещение, --- главная методологическая претензия к текущей практике.

\subsection{Условное правдоподобие схемы A}\label{sec:condlik}

Предложение~\ref{prop:gamma} даёт больше, чем разбор смещения. Полагая $t_0=\ln n_{\min}$
и приняв $(N_b,D_b)$ за данные по основанию~$b$, получаем логарифмическое правдоподобие
\begin{equation}\label{eq:loglik}
\ell_{\mathrm A}(\kappa)=\sum_b\Bigl[N_b\ln\frac{\kappa}{\ln b}-\frac{\kappa}{\ln b}D_b\Bigr]+\mathrm{const},
\end{equation}
распределение которого не зависит от границы поиска $L_b$. Это существенно практически:
$L_b$ документирована лишь при $b=2$ и $b=10$. Дифференцируя~\eqref{eq:loglik} и полагая
$S=\sum_b D_b/\ln b$, получаем
\begin{equation}\label{eq:mle}
\hat\kappa=\frac{M}{S},\qquad \tilde\kappa=\frac{M-1}{S},\qquad M=\sum_b N_b .
\end{equation}
Так как $\kappa S=\sum_b\lambda_bD_b$ есть сумма независимых величин $\Gam(N_b,1)$, то
\begin{equation}\label{eq:pivot}
\kappa S\sim\Gam(M,1),
\end{equation}
то есть точная центральная статистика; доверительные интервалы и $p$-значения получаются
в замкнутой форме. \emph{Ровно это утверждение и требует ревизии} --- см.~\S\ref{sec:scheme}.

\begin{remark}[Независимость по основаниям]\label{rem:indep}
Объединение предполагает независимость процессов по разным основаниям. Строго она не
выполняется на малых индексах: семнадцать значений индекса встречаются более чем в одном
основании, причём $n=3$ --- в девяти, поскольку $\Rb(3)=b^2+b+1$ просто для многих~$b$.
Все такие совпадения лежат при $n\le317$. Существеннее, впрочем, не совпадение индексов
(числа $\Rb(n)$ при разных $b$ различны), а возможная общность алгебраических делителей;
и то и другое затрагивает лишь малые~$n$, тогда как $D_b$ определяется верхним концом
диапазона. Раздел~\ref{sec:robust} подтверждает, что отбрасывание всех индексов ниже~$10^3$
вывод не меняет. Отметим, что к условному мультиномиальному критерию
\S\ref{sec:homog} это замечание относится сильнее, чем к объединённой оценке:
зависимость искажает не среднее, а нулевое распределение статистики; та же
проверка с усечением покрывает и его.
\end{remark}

\section{Две схемы наблюдения}\label{sec:scheme}

\subsection{В чём различие}

Пивот~\eqref{eq:pivot} выведен в предположении, что $N_b$ задано планом наблюдения:
мы смотрим на процесс, пока не наберём $N_b$ событий, и останавливаемся. Назовём это
\emph{схемой~A} (обратная выборка, цензурирование типа~II).

Данные получены иначе. Для каждого основания вычислительный поиск покрыл интервал
$(t_0,\ln L_b]$ целиком, и найденными оказались все события, попавшие внутрь; число
находок $N_b$ --- случайная величина, а фиксирован фронт $L_b$. Назовём это
\emph{схемой~B}.

\begin{remark}[Что здесь классика, а что нет]\label{rem:priority}
Само различие схем классическое и восходит к теории надёжности, где схема~A
называется \emph{failure-truncated}, а схема~B --- \emph{time-truncated}
испытанием на долговечность~\cite{epstein1954trunc,epsteinsobel1954}. Оценка
$(N-1)/D$, несмещённая при схеме~A (следствие~\ref{cor:kappa}), там
учебниковая, и мы не претендуем на неё. Распределение ОМП среднего при
цензурировании типа~I изучено Бартоломью~\cite{bartholomew1963}. Отметим также
элементарный факт: при схеме~B с известным окном достаточной статистикой служит
$N\sim\Pois(\lambda W)$, а $\mathbb E f(N)=\sum_nf(n)e^{-\mu}\mu^n/n!$ есть целая
функция от $\mu=\lambda W$, которая не может совпадать с $W/\mu$, имеющей полюс в
нуле. Поэтому несмещённой оценки \emph{среднего} $1/\lambda$ при схеме~B не
существует (для цензурированных по времени выборок из экспоненциального
распределения см.~\cite{sengupta2009}); то же рассуждение применимо и условно на
$D_b$. Предложение~\ref{prop:cond}
этому не противоречит: оно даёт точную несмещённость оценки \emph{интенсивности}.

Новым мы считаем другое. В классической постановке момент прекращения
испытания \emph{известен}: время $T$ назначает экспериментатор. Здесь фронт
$L_b$ документирован лишь для $b=2$ и $b=10$ (\S\ref{sec:frontiers}), поэтому полную
экспозицию приходится не брать из протокола, а восстанавливать --- через
отсутствие последействия, из самих находок. Отсюда и берётся поправка
\emph{на каждое основание}, а не на выборку в целом: неизвестных фронтов
столько же, сколько оснований. Обзор смещений, порождаемых краем окна
наблюдения, но в непараметрической постановке и для формы распределения
интервалов, а не для интенсивности, дан в~\cite{kivela2015}.

Двойственные задачи известны давно. В палеонтологии по находкам окаменелостей
оценивают неизвестный конец стратиграфического интервала, и наблюдённый крайний
уровень продлевают примерно на один средний промежуток между
находками~\cite{straussSadler1989,marshall1990}. В экологии по записям
наблюдений вида оценивают неизвестный момент его вымирания~\cite{solow1993}.
Общая задача об оценке точки усечения рассмотрена в~\cite{robsonwhitlock1964}.
Там, однако, неизвестный конец окна --- параметр интереса. У нас он мешающий
параметр, и оценивается интенсивность. Ближе всего по технике
работа~\cite{liu2023}: там окно наблюдения начинается с первого события, и
правдоподобие факторизуется на пуассоновский множитель и множитель, не зависящий
от интенсивности. Работ, где цензурирование типа~I с незаписанным моментом
цензурирования рассматривалось бы как задача об объединённой оценке
интенсивности по многим стратам, мы не нашли.
\end{remark}

Различие не косметическое. При схеме~B, условно на $N_b=n$, положения событий суть
$n$ независимых равномерных величин на $(t_0,\ln L_b]$, поэтому условное распределение
$D_b$ зависит от $L_b$ и \emph{не} зависит от $\lambda_b$: при фиксированном $N_b$
величина $D_b$ информативна о фронте, а не об интенсивности.

Точнее, при схеме~B с известной шириной окна $W_b=\ln L_b-t_0$ правдоподобие
основания~$b$ есть $\lambda_b^{N_b}e^{-\lambda_bW_b}$, а при неизвестной оно
зависит от мешающего параметра $W_b\ge D_b$ и монотонно убывает по нему. Совместная
ОМП поэтому ставит $\hat W_b=D_b$ --- фронт ровно на последнюю находку, --- и
выражение~\eqref{eq:loglik} оказывается \emph{профильным} правдоподобием по
мешающим параметрам $W_1,\dots,W_B$, а $M/S$ --- профильной ОМП. Число мешающих
параметров растёт вместе с числом оснований, и профильная ОМП несостоятельна:
при $B\to\infty$ и ограниченных $\lambda_bW_b$ она сходится не к $\kappa$. Это
классическая ситуация Неймана --- Скотта~\cite{neymanscott1948}
(обзор --- \cite{lancaster2000}), и стандартный выход из неё --- обусловливание
на статистике, достаточной для мешающих параметров
\cite{kalbfleischsprott1970,andersen1970}. Его и реализует \S\ref{sec:exact}.

Заметим также, что \S\ref{sec:sameerror} настоящей работы, где для $b=2$ фронт известен
и используется явно, построен именно на схеме~B. Обе схемы в одной работе применять
нельзя; ниже мы принимаем схему~B как соответствующую способу получения данных
и строим для неё точный вывод.

\subsection{Точный условный закон схемы B}\label{sec:exact}

Ключевое наблюдение состоит в том, что при схеме~B обусловливание надо проводить
не на числе событий, а на положении последнего из них. Тогда неизвестный фронт
исчезает из задачи полностью. Требуется одно допущение.

\medskip
\noindent\textbf{(A1)} \emph{Положение фронта $L_b$ не выбиралось по результатам поиска:
вычислительный бюджет, а не находки, определяет, до какого $n$ доведено просеивание.
Формально: ширина окна $W_b=\ln L_b-t_0$ есть случайная величина с произвольным
неизвестным распределением $\pi_b$ (в частности, константа), не зависящая от
процесса находок по основанию~$b$.}

\medskip
\noindent На языке анализа выживаемости (A1) есть требование
\emph{неинформативности} цензурирования: момент остановки не зависит от того,
что процесс успел породить. Распределения $\pi_b$ могут быть любыми и различаться
по основаниям; они нигде не оцениваются.

\begin{proposition}[Точное условное правдоподобие]\label{prop:cond}
Пусть при схеме~B однородный пуассоновский процесс интенсивности $\lambda_b$
наблюдается на $(t_0,t_0+W_b]$, выполнено (A1), $N_b\ge1$ и $D_b=t_{N_b}-t_0$ ---
расстояние от начала отсчёта до последнего найденного события.
\begin{enumerate}
\item[\rm(i)] Совместная плотность $(N_b,D_b)$ в точке $(n,d)$, $n\ge1$, $d>0$,
факторизуется:
\begin{equation}\label{eq:factor}
p(n,d)=\underbrace{e^{-\lambda_bd}\frac{(\lambda_bd)^{n-1}}{(n-1)!}}_{\text{зависит от }\lambda_b}
\;\cdot\;
\underbrace{\lambda_b\int_{[d,\infty)}e^{-\lambda_b(w-d)}\,\pi_b(dw)}_{\text{плотность }D_b},
\end{equation}
так что при фиксированной $\lambda_b$ статистика $D_b$ достаточна для $W_b$
(и для всего распределения $\pi_b$).
\item[\rm(ii)] Следовательно,
\begin{equation}\label{eq:condlaw}
K_b:=N_b-1\ \bigm|\ D_b\ \sim\ \Pois(\lambda_b D_b),
\end{equation}
и это распределение не зависит ни от $W_b$, ни от $\pi_b$.
\item[\rm(iii)] Для независимых по основаниям процессов
\begin{equation}\label{eq:pivotB}
K:=M-B\ \bigm|\ (D_b)_{b}\ \sim\ \Pois(\kappa S),\qquad
S=\sum_b\frac{D_b}{\ln b},
\end{equation}
и $\hat\kappa_{\mathrm B}=(M-B)/S$ есть оценка условного максимального
правдоподобия, точно условно несмещённая:
$\mathbb E[\hat\kappa_{\mathrm B}\mid (D_b)_b]=\kappa$.
\end{enumerate}
\end{proposition}

\begin{proof}
(i) Зафиксируем сначала $W_b=w$. Событие $\{N_b=n,\ D_b\in[d,d+\Delta)\}$ при
$0<d<d+\Delta\le w$ и $\Delta\to0$ означает: на $(t_0,t_0+d)$ ровно $n-1$
событий, в $[t_0+d,t_0+d+\Delta)$ одно, на $[t_0+d+\Delta,t_0+w]$ ни одного. По
независимости приращений пуассоновского процесса его вероятность равна
$e^{-\lambda_bd}\frac{(\lambda_bd)^{n-1}}{(n-1)!}\cdot\lambda_b\Delta\cdot
e^{-\lambda_b(w-d)}+o(\Delta)$, а при $d>w$ равна нулю. По (A1) усреднение по
$w\sim\pi_b$ затрагивает лишь последний множитель и даёт~\eqref{eq:factor}.
Суммирование по $n\ge1$ показывает, что второй множитель в~\eqref{eq:factor} есть
маргинальная плотность $D_b$.
(ii) Деление~\eqref{eq:factor} на маргинальную плотность $D_b$.
(iii) Сумма независимых пуассоновских величин пуассоновская; логарифм условного
правдоподобия $\ell_{\mathrm B}(\kappa)=\sum_b\bigl[K_b\ln\kappa-\kappa D_b/\ln b\bigr]
+\mathrm{const}$ максимизируется при $\kappa=K/S$, а условная несмещённость
немедленна из $\mathbb E[K\mid D]=\kappa S$.
\end{proof}

\begin{remark}[Условное и интегрированное правдоподобие]\label{rem:integrated}
Тот же результат получается интегрированием правдоподобия по мешающему
параметру. При фиксированном окне совместная плотность пропорциональна
$\lambda_b^{n}d^{\,n-1}e^{-\lambda_bW_b}\mathbf 1_{\{W_b\ge d\}}$, и
$\int_d^\infty\lambda_b^{n}e^{-\lambda_bW_b}\,dW_b=\lambda_b^{n-1}e^{-\lambda_bd}$ ---
ровно ядро~\eqref{eq:condlaw}. Это интегрированное правдоподобие в
смысле~\cite{berger1999} с несобственной равномерной мерой на $W_b$. Такая
процедура сама по себе зависит от выбора меры: мера, равномерная по
$\mu_b=\lambda_bW_b$, то есть $d\mu_b=\lambda_b\,dW_b$, дала бы
$\lambda_b^{n}e^{-\lambda_bd}$ и вернула бы несостоятельную оценку $M/S$. Здесь, однако,
выбор не произволен. Во-первых, результат совпадает с \emph{условным}
правдоподобием~\eqref{eq:condlaw}, которое определено без всякой меры и
инвариантно относительно перепараметризации $W_b$. Во-вторых, $W_b$ входит
в~\eqref{eq:factor} как параметр сдвига величины $W_b-D_b$, и $dW_b$ ---
инвариантная мера группы сдвигов. В-третьих, при $\pi_b(dw)=dw$ маргинальный
множитель в~\eqref{eq:factor} равен $\lambda_b\int_d^\infty e^{-\lambda_b(w-d)}dw=1$
и от $\lambda_b$ не зависит. Значит, отбрасываемый при обусловливании множитель
при неизвестном фронте не несёт информации о $\kappa$, которую можно было бы
использовать, не задавая $\pi_b$.
\end{remark}

\begin{remark}[Связь с задачей Неймана --- Скотта]\label{rem:neymanscott}
Поправка <<одно событие на основание>> имеет ту же природу, что потеря одной
степени свободы на группу в классическом примере Неймана ---
Скотта~\cite{neymanscott1948}. При $y_{ij}\sim N(\mu_i,\sigma^2)$ профильная ОМП
$\hat\sigma^2$ несостоятельна, а обусловливание на $\bar y_i$, достаточной для
$\mu_i$ при фиксированной $\sigma^2$, даёт
$\sum_j(y_{ij}-\bar y_i)^2\sim\sigma^2\chi^2_{n_i-1}$. В обоих случаях мешающий
параметр группы одномерен и является параметром сдвига, а достаточная для него
статистика <<поглощает>> ровно одно наблюдение: среднее --- одну степень свободы,
последнее событие --- одно событие. Общий механизм --- условное правдоподобие при
статистике, достаточной для мешающего параметра при фиксированном параметре
интереса~\cite{kalbfleischsprott1970}. Его состоятельность при растущем числе
страт в регулярных условиях установлена Андерсеном~\cite{andersen1970}, а
оптимальность условной оценки в моделях смеси по мешающему параметру ---
Линдсеем~\cite{lindsay1980}.
\end{remark}

\begin{remark}[Неоднородная интенсивность]\label{rem:inhom}
Рассуждение~(i) не использует однородности. Если интенсивность
$\lambda_b(t;\theta)$ зависит от параметра $\theta$, то условно на $D_b$ события
на $(t_0,t_0+D_b)$ образуют пуассоновский процесс той же интенсивности, и
условный логарифм правдоподобия есть
\begin{equation}\label{eq:condinhom}
\ell_{\mathrm B}(\theta)=\sum_b\Bigl[\sum_{i=1}^{N_b-1}\ln\lambda_b(t_i;\theta)
-\int_{t_0}^{t_0+D_b}\lambda_b(t;\theta)\,dt\Bigr]:
\end{equation}
последнее событие исключается из суммы, а интеграл обрывается на нём. Именно в
таком виде следует оценивать модели \S\ref{sec:struct} с переменной
интенсивностью.
\end{remark}

Практические следствия предложения~\ref{prop:cond} стоит выписать отдельно,
поскольку они устраняют почти все оговорки предыдущей редакции этой работы.

\begin{enumerate}
\item \emph{Точные интервалы без симуляции.} Инверсия пуассоновской ставки даёт
интервал Гарвуда~\cite{garwood1936}
\begin{equation}\label{eq:garwood}
\Bigl[\tfrac{1}{2S}\chi^2_{\alpha/2,\,2K},\ \tfrac{1}{2S}\chi^2_{1-\alpha/2,\,2(K+1)}\Bigr],
\end{equation}
покрытие которого никогда не ниже номинала. При нашем $S$ оно равно $0{,}9540$
при $\kappa=\eg$ и лежит в пределах $0{,}950$--$0{,}958$ при
$\kappa\in[1{,}5;\,2{,}2]$ (превышение --- обычный эффект дискретности).
\item \emph{Никакого остатка.} Ни экспоненциально малых членов, ни оговорок об
основаниях без находок: обусловливание на $N_b\ge1$ выполняется по построению,
поскольку $D_b$ определено только тогда. Основание без единой находки в
условное правдоподобие не входит.
\item \emph{Никакой зависимости от положения фронта.} Правило, по которому
выбирался $L_b$, в~\eqref{eq:condlaw} не входит вовсе --- при условии (A1).
\item \emph{Цена обусловливания мала.} $\operatorname{SD}(\hat\kappa_{\mathrm B})
=\sqrt{\kappa/S}=0{,}1270$ против $\sqrt{\kappa/S_{\mathrm{full}}}=0{,}1224$
при известных фронтах, стоящих на своём ожидаемом месте,
$S_{\mathrm{full}}\approx S+B/\kappa$ (предложение~\ref{prop:delta}). Отказ от
знания $L_b$ стоит $3{,}6\%$ стандартного отклонения, или $7{,}1\%$ фишеровской
информации: $B/\kappa$ из $S+B/\kappa$, то есть примерно $B$ событий из $M$.
Знание фронтов нужно не столько для точности, сколько для проверки (A1)
(\S\ref{sec:noescape}).
\end{enumerate}

\noindent При нарушении (A1) обусловливание на $D_b$ перестаёт быть безобидным.
Если поиск прекращается ровно на находке, то $D_b$ и есть момент остановки,
$\delta_b=0$, и возвращается схема~A. Если же поиск прекращается после
достаточно долгой полосы без находок, то сама остановка несёт информацию о
процессе до $D_b$, и закон~\eqref{eq:condlaw} перестаёт быть верным
(замечание~\ref{rem:bracket}).

\subsection{Сколько стоит незнание фронта}\label{sec:cost}

Помимо замечания~\ref{rem:neymanscott}, величину поправки поясняет прямой
моментный расчёт. Мы сохраняем его как содержательное толкование и как основу
анализа нарушений (A1) в замечании~\ref{rem:bracket}.

\begin{proposition}\label{prop:delta}
Пусть при схеме~B поиск по основанию~$b$ покрыл $(t_0,\ln L_b]$, ширина окна
$W_b=\ln L_b-t_0$ фиксирована (при (A1) --- условно на $W_b$), а
$\delta_b=\ln L_b-t_{N_b}$ --- расстояние от фронта до последнего
найденного события (при $N_b=0$ полагаем $\delta_b=W_b$). Тогда
\begin{equation}\label{eq:delta}
\Pr(\delta_b>x)=e^{-\lambda_b x}\ \ (0\le x<W_b),\qquad
\Pr(\delta_b=W_b)=e^{-\lambda_b W_b},
\end{equation}
то есть $\delta_b$ есть экспоненциальная величина с параметром $\lambda_b$,
усечённая на $W_b$, и
\begin{equation}\label{eq:Edelta}
\mathbb E[\delta_b]=\frac{1-e^{-\lambda_b W_b}}{\lambda_b}\ \le\ \frac{1}{\lambda_b}.
\end{equation}
Полная экспозиция есть $S_{\mathrm{full}}=S+\sum_b\delta_b/\ln b=\sum_b W_b/\ln b$,
причём $\mathbb E[M]=\kappa S_{\mathrm{full}}$ тождественно, а
\begin{equation}\label{eq:EBd}
\mathbb E\Bigl[\kappa\sum_b\delta_b/\ln b\Bigr]=B-\sum_b e^{-\lambda_b W_b}.
\end{equation}
\end{proposition}

\begin{proof}
Событие $\{\delta_b>x\}$ при $x<W_b$ означает, что на интервале
$(\ln L_b-x,\ln L_b]$ длины $x$ нет ни одного события процесса; вероятность этого
равна $e^{-\lambda_b x}$. При $x\to W_b^-$ остаётся случай <<событий нет вовсе>>,
что и даёт атом в~\eqref{eq:delta}. Интегрируя хвост,
$\mathbb E[\delta_b]=\int_0^{W_b}e^{-\lambda_b x}\,dx=(1-e^{-\lambda_b W_b})/\lambda_b$,
откуда~\eqref{eq:Edelta}, и неравенство очевидно. Далее, $S_{\mathrm{full}}$ не
случайна: по построению $D_b+\delta_b=W_b$, поэтому
$\kappa S_{\mathrm{full}}=\sum_b\lambda_b W_b=\mathbb E[M]$ --- это в точности
уравнение моментов. Наконец,
$\kappa\,\mathbb E[\delta_b]/\ln b=\lambda_b\mathbb E[\delta_b]=1-e^{-\lambda_b W_b}$,
и суммирование по $b$ даёт~\eqref{eq:EBd}.
\end{proof}

Иными словами, незнание фронта стоит информации, эквивалентной одному событию
\emph{на каждое основание}, --- а формула~\eqref{eq:mle} вычитает одно событие
\emph{всего}. При $B=15$ разница составляет четырнадцать событий, то есть $6{,}9\%$:
\begin{equation}\label{eq:cost}
\frac{M-1}{S}=1{,}9648\qquad\text{против}\qquad\frac{M-B}{S}=1{,}8381 .
\end{equation}
Отбрасываемый в~\eqref{eq:EBd} остаток $\sum_b e^{-\lambda_b W_b}$ есть ожидаемое
число оснований, у которых поиск не дал ни одной находки; в наших данных
$\lambda_bW_b\approx N_b\ge7$, так что он не превосходит $3\cdot10^{-3}$ против
$B=15$. Подчеркнём, что в точной формулировке~\eqref{eq:condlaw} этот остаток не
возникает вовсе: моментный расчёт нужен здесь только для интерпретации.

\begin{remark}[Что остаётся без (A1)]\label{rem:bracket}
В предыдущей редакции утверждалось, что семейство $(M-Bm)/S$ с $m\in[0,1]$ даёт
доказуемую двустороннюю вилку $[(M-B)/S,\ (M-1)/S]$. Это утверждение
\emph{неверно}, и мы его отзываем. Неравенство~\eqref{eq:Edelta} выведено при
фиксированном (или независимом от процесса) окне, то есть именно при (A1), а
тогда вилка не нужна: верна оценка $(M-B)/S$. Если же (A1) нарушено, то
$\lambda_b\mathbb E[\delta_b]$ может быть любым неотрицательным числом.

Строго сохраняется лишь \emph{верхняя} граница. Если фронт есть момент
остановки относительно истории поиска (решение прекратить поиск не заглядывает
вперёд), то по тождеству Вальда для пуассоновского процесса
$\mathbb E[N_b]=\lambda_b\,\mathbb E[W_b]$ \emph{при любом} правиле остановки с
$\mathbb E[W_b]<\infty$. Поскольку $W_b=D_b+\delta_b\ge D_b$, получаем
$\mathbb E[M]\ge\kappa\,\mathbb E[S]$: в смысле моментов $M/S$ ограничивает $\kappa$
сверху. Нижней границы без допущений нет. Правило <<прекратить поиск после полосы
длиной $g$ средних интервалов без находок>> даёт $\lambda_b\delta_b\ge g$. Кроме
того, само ожидание такой полосы уплотняет события перед $D_b$, так что
закон~\eqref{eq:condlaw} перестаёт быть верным. Таблица~\ref{tab:stop} показывает
последствия: при $g=2$ оценка $(M-B)/S$ завышена на $6{,}7\%$, а точный критерий
предложения~\ref{prop:cond} отвергает верную гипотезу в $18\%$ случаев вместо
$5\%$. Правило <<остановиться на следующей находке>> ($\delta_b=0$) смещает
$(M-B)/S$ вниз на $6{,}4\%$, и тогда правильна оценка $(M-1)/S$.

Поэтому точная формулировка такова. При (A1) основной результат~\eqref{eq:resB}
точен. При остановке <<на находке>> правильна оценка~\eqref{eq:resA}, и она же
в смысле моментов ограничивает $\kappa$ сверху при любом правиле. При остановке
<<после засухи>> истинное $\kappa$ \emph{ниже} $(M-B)/S$, то есть ещё ближе к
$\eg$ или ниже его. Ни одно из нарушений (A1) не создаёт превышения над ЛПВ,
которое данные требовали бы объяснять. Но величина неопределённости без
документированных фронтов не ограничена снизу: это довод в пользу
публикации $L_b$ сильнее того, что приводился раньше.
\end{remark}

\begin{table}[t]
\centering
\caption{Схема~B при различных правилах остановки поиска: относительное среднее
двух объединённых оценок, отношение $\mathbb E[M]/(\kappa\mathbb E[S])$ (по тождеству
Вальда $\ge1$ при любом правиле) и ошибка I рода точного двустороннего
пуассоновского критерия~\eqref{eq:pivotB} на уровне $5\%$. Бюджет
$W^0_b=N_b/\lambda_b$; <<засуха $g$>>: после исчерпания бюджета поиск
продолжается, пока не пройдёт $g/\lambda_b$ без находок. $20\,000$ реализаций на
строку, точность около $\pm0{,}05\%$ по средним и $\pm0{,}003$ по уровню (скрипт
\texttt{stopping\_rules.py}, тег~[S13], зерно $20260924$).
\encaptab{Scheme~B under various stopping rules: relative mean of the two pooled
estimators, the ratio $\mathbb E[M]/(\kappa\mathbb E[S])$ (at least $1$ under every
rule by Wald's identity), and the type~I error of the exact two-sided Poisson
test~\eqref{eq:pivotB} at the $5\%$ level. Budget $W^0_b=N_b/\lambda_b$;
``drought $g$'': after the budget is spent the search continues until
$g/\lambda_b$ passes without a find. $20\,000$ replications per row, accuracy
about $\pm0.05\%$ for the means and $\pm0.003$ for the level (script
\texttt{stopping\_rules.py}, tag~[S13], seed $20260924$).}}
\label{tab:stop}
\small
\begin{tabular}{lrrrr}
\toprule
Правило остановки & $\mathbb E[(M{-}1)/S]/\kappa$ & $\mathbb E[(M{-}B)/S]/\kappa$
& $\dfrac{\mathbb E[M]}{\kappa\mathbb E[S]}$ & ошибка I рода\\
\midrule
(A1): фронт $=W^0_b$ & 1{,}0688 & 0{,}9999 & 1{,}0737 & 0{,}046\\
(A1): фронт $=UW^0_b$, $U\sim\mathrm U(0{,}5;\,2)$ & 1{,}0551 & 1{,}0000 & 1{,}0582 & 0{,}043\\
остановка на находке & 0{,}9963 & 0{,}9362 & 1{,}0003 & 0{,}165\\
засуха $g=0{,}5$ & 1{,}0762 & 1{,}0074 & 1{,}0811 & 0{,}050\\
засуха $g=1$ & 1{,}0923 & 1{,}0250 & 1{,}0973 & 0{,}063\\
засуха $g=2$ & 1{,}1227 & 1{,}0667 & 1{,}1277 & 0{,}178\\
\bottomrule
\end{tabular}
\end{table}

\subsection{Калибровка старого пивота}\label{sec:calib}

Насколько плох гамма-пивот~\eqref{eq:pivot}, если данные порождены схемой~B,
зависит от того, как далеко ушёл фронт, --- и это стоит приводить диапазоном, а
не одним числом. Таблица~\ref{tab:calib} даёт обе величины сразу: ошибку I рода
старого критерия и относительное смещение обеих объединённых оценок.

\begin{table}[t]
\centering
\caption{Схема~B при семи правилах расположения фронта: ошибка I рода
номинального $5\%$ критерия, построенного на гамма-пивоте~\eqref{eq:pivot}, и
относительное смещение двух объединённых оценок. Правила заданы через ширину окна
$W_b=\ln L_b-t_0$: <<$\mathbb E[\#]$ = наблюдённое>> есть $W_b=N_b/\lambda_b$;
<<$+s$ средних интервалов>> --- $W_b=(N_b+s)/\lambda_b$; <<в $c$ раз дальше по $t$>>
--- $W_b=c\,D_b$ с наблюдённым $D_b$. $50\,000$ реализаций на строку, точность
около $\pm0{,}003$ по уровню и $\pm0{,}1\%$ по средним (тег~[S10]).
\encaptab{Scheme~B under seven frontier rules: type~I error of the nominal $5\%$
test built on the Gamma pivot~\eqref{eq:pivot}, and the relative bias of the two
pooled estimators. The rules are stated through the window width
$W_b=\ln L_b-t_0$: ``$\mathbb E[\#]$ = observed'' is $W_b=N_b/\lambda_b$;
``$+s$ mean intervals'' is $W_b=(N_b+s)/\lambda_b$; ``$c$ times further in $t$''
is $W_b=c\,D_b$ with the observed $D_b$. $50\,000$ replications per row, accuracy
about $\pm0.003$ for the level and $\pm0.1\%$ for the means (tag~[S10]).}}
\label{tab:calib}
\small
\begin{tabular}{lrrr}
\toprule
Правило для фронта & ошибка I рода & $\mathbb E[(M{-}1)/S]/\kappa$
& $\mathbb E[(M{-}B)/S]/\kappa$\\
\midrule
$\mathbb E[\#\text{событий}]$ = наблюдённое & 0{,}1698 & 1{,}0696 & 1{,}0006\\
фронт $+0{,}5$ среднего интервала & 0{,}1592 & 1{,}0658 & 0{,}9993\\
фронт $+1$ средний интервал & 0{,}1617 & 1{,}0652 & 1{,}0009\\
фронт $+2$ средних интервала & 0{,}1512 & 1{,}0606 & 1{,}0005\\
фронт $+5$ средних интервалов & 0{,}1356 & 1{,}0505 & 1{,}0001\\
фронт в $1{,}5$ раза дальше по $t$ & 0{,}1335 & 1{,}0497 & 0{,}9997\\
фронт вдвое дальше по $t$ & 0{,}1123 & 1{,}0369 & 0{,}9999\\
\midrule
A (контроль) & 0{,}0494 & --- & ---\\
\bottomrule
\end{tabular}
\end{table}

Три вывода из таблицы~\ref{tab:calib}. Во-первых, номинальный $5\%$ критерий при
схеме~B отвергает верную гипотезу в $11$--$17\%$ случаев: заявленная точность
интервалов завышена вдвое-втрое, причём тем сильнее, чем ближе фронт к последней
находке. Во-вторых, оценка $(M-1)/S$ смещена вверх на $4$--$7\%$ при всех правилах,
тогда как $(M-B)/S$ несмещена с точностью до шума Монте-Карло при всех правилах ---
что и предсказывает предложение~\ref{prop:cond}, поскольку условная несмещённость
влечёт безусловную. В-третьих, само число <<$16{,}2\%$>>, фигурировавшее в более
ранней редакции, есть значение при одном конкретном сценарии; корректно приводить
диапазон. Все семь правил таблицы~\ref{tab:calib} задают фронт без оглядки на
находки, то есть удовлетворяют (A1); правила, нарушающие (A1), разобраны в
замечании~\ref{rem:bracket} и таблице~\ref{tab:stop}.

\subsection{Документированные фронты: эмпирическая проверка (A1)}\label{sec:frontiers}

Датированный и атрибутированный фронт выше последней находки удалось найти для
четырёх оснований (проверено 24.09.2026):
\begin{itemize}
\item $b=2$: отчёт GIMPS на момент считывания данных, $X_{\mathrm{ver}}=81\,648\,221$
и $X_{\mathrm{ft}}=141\,308\,443$ (\S\ref{sec:data});
\item $b=10$: <<Search limit is 10800000, currently>> (С.~Баталов, A004023,
1.07.2021);
\item $b=22$ и $b=26$: <<No other terms < 100000>> (Р.~Прайс, A127997, A127999,
2012~г.; последние находки $27\,823$ и $26\,717$ сделаны в 2007--2008~гг.).
\end{itemize}
Остальные упоминания пределов поиска либо устарели и лежат далеко ниже $n_{\max}$,
либо совпадают с последней находкой и отмечают её, а не поиск.

При (A1) величины $\lambda_b\delta_b$ независимы и распределены как $\mathrm{Exp}(1)$
(усечение на $W_b$ здесь пренебрежимо), так что их сумма по $k$ основаниям имеет
распределение $\Gam(k,1)$. Малые значения указывают на остановку <<на находке>>
($m<1$), большие --- на остановку <<после засухи>> ($m>1$,
замечание~\ref{rem:bracket}). Наблюдаются значения $0{,}143$ ($b=2$ по
$X_{\mathrm{ver}}$; $0{,}093$ по $X_{\mathrm{ft}}$), $0{,}215$ ($b=10$), $0{,}737$
($b=22$) и $0{,}722$ ($b=26$). Среднее равно $0{,}44$--$0{,}45$ при ожидаемой единице.
Вероятность столь малой суммы при (A1) равна $0{,}10$--$0{,}11$ (тег~[S18]).

Выводы осторожны: наблюдений четыре, фронт $b=2$ движется непрерывно, а два
основания входят лишь в расширенный набор. Тем не менее это первая проверка (A1)
по данным, и её результат информативен по \emph{направлению}. Документированные
фронты расположены не дальше, а ближе к последним находкам, чем предсказывает (A1).
Опасное для основного вывода нарушение --- остановка <<после засухи>>, смещающая
$(M-B)/S$ вверх, --- данными не поддерживается: верхний хвост равен $0{,}89$--$0{,}90$.
Если (A1) и нарушено, то в сторону $m<1$, и истинное значение лежит между
$(M-B)/S$ и $(M-1)/S$.

Документированные фронты можно и использовать напрямую. Для оснований с известным
фронтом применимо полное правдоподобие: $N_b\sim\Pois(\lambda_bW_b)$ со всеми
событиями и полным окном. Для остальных применим условный закон~\eqref{eq:condlaw}.
Сумма снова точно пуассоновская. Гибридная оценка с известными фронтами $b=2$
($X_{\mathrm{ft}}$) и $b=10$ равна $\hat\kappa=1{,}853$ при точном $95\%$
интервале $[1{,}608;\,2{,}125]$ и $p=0{,}588$ против $\eg$. Она отличается от
$(M-B)/S=1{,}838$ на $0{,}8\%$, и основной вывод не меняется.

\subsection{Данные не выбирают схему за нас}\label{sec:noescape}

Подчеркнём, что $(M-1)/S$ и $(M-B)/S$ --- функции одних и тех же достаточных статистик
$(N_b,D_b)$. Никакая обработка имеющихся данных не скажет, какая из них верна: выбор
определяется тем, как данные собирались, и цена его составляет $6{,}9\%$ --- вдвое
больше, чем весь остаточный разрыв с ЛПВ ($3{,}2\%$).

Естественная попытка обойти выбор --- строить объединённую оценку из заведомо
несмещённых кусков.
Оценка по одному основанию $\tilde\kappa_b=(N_b-1)\ln b/D_b$ несмещена при схеме~A
(следствие~\ref{cor:kappa}) и несмещена при схеме~B с точностью до экспоненциально
малого члена. Действительно, при $N_b=n\ge2$ отношение $D_b/W_b$ распределено как
$\mathrm{Beta}(n,1)$, и тождество
$\mathbb E[1/\mathrm{Beta}(\alpha,\beta)]=(\alpha+\beta-1)/(\alpha-1)$ даёт
$\mathbb E[\tilde\lambda_b\mid N_b=n]=n/W_b$; при $n\in\{0,1\}$ оценка равна нулю.
Усреднение по $n\sim\mathrm{Pois}(\lambda_bW_b)$ поэтому даёт не $\lambda_b$, а
\[
  \mathbb E[\tilde\lambda_b]=\lambda_b-\frac{\Pr(N_b=1)}{W_b}
  =\lambda_b\Bigl(1-e^{-\lambda_bW_b}\Bigr),
\]
то есть тот же экспоненциально малый дефицит: при $\lambda_bW_b\approx N_b\ge7$ он
не превосходит $0{,}1\%$.

Тем не менее взвешенное среднее
$$\kappa_\ast=\sum\nolimits_b(N_b-2)\tilde\kappa_b\Big/\sum\nolimits_b(N_b-2)$$
с обратно-дисперсионными весами даёт при схеме~A смещение $+0{,}14\%$,
а при схеме~B --- $+7{,}5\%$, то есть практически то же, что и $(M-1)/S$. Причина в том,
что веса --- сами числа событий --- при схеме~B случайны и положительно коррелируют с
оценками; несмещённость каждого слагаемого не переносится на отношение сумм.
Правильный способ объединения --- не взвешивать по-основанные оценки, а
суммировать достаточные статистики, как это делает~\eqref{eq:pivotB}.

Отсюда мы делаем узкий, но, как нам кажется, важный вывод: \emph{выигрыш в точности от
объединения оснований неотделим от предположения о схеме наблюдения}.
Единственный радикальный выход --- документировать фронты поиска $L_b$, как это делает
GIMPS для $b=2$. Мы обращаемся с этим к участникам соответствующих поисковых проектов.

\section{Результаты}\label{sec:results}

\subsection{Оценки по основаниям}

\begin{table}[t]
\centering
\caption{Оценки по основаниям. $\hat G_b=D_b/(N_b\ln b)$ --- концевая форма;
$\hat G_b^{\mathrm{OLS}}$ --- наклон МНК $\ln n_k$ на $k$, делённый на $\ln b$; обе
несмещены (предложение~\ref{prop:unbiased}). $\hat\kappa_b=K_b/(D_b/\ln b)$ ---
условная оценка схемы~B с точным интервалом Гарвуда~\eqref{eq:garwood} из
$\Pois(\kappa D_b/\ln b)$. Последний столбец --- mid-PIT пуассоновского
распределения при $\kappa=\eg$: значения вблизи единицы отвечают основаниям с
\emph{высокой} оценкой $\kappa_b$ (ориентация обратна гамма-PIT прежней
редакции). ЛПВ предсказывает $G=\emg=0{,}5615$ и $\kappa=\eg=1{,}7811$.
\encaptab{Per-base estimates. $\hat G_b=D_b/(N_b\ln b)$ is the endpoint form;
$\hat G_b^{\mathrm{OLS}}$ is the OLS slope of $\ln n_k$ on $k$ divided by $\ln b$;
both are unbiased (Proposition~\ref{prop:unbiased}).
$\hat\kappa_b=K_b/(D_b/\ln b)$ is the scheme~B conditional estimator with the
exact Garwood interval~\eqref{eq:garwood} from $\Pois(\kappa D_b/\ln b)$. The last
column is the Poisson mid-PIT at $\kappa=\eg$: values near one correspond to bases
with a \emph{high} $\kappa_b$ (orientation opposite to the Gamma PIT of the earlier
version). LPW predicts $G=\emg=0.5615$ and $\kappa=\eg=1.7811$.}}
\label{tab:bases}
\small
\begin{tabular}{rrrrrrrl r}
\toprule
$b$ & $N_b$ & $K_b$ & $\hat G_b$ & $\hat G_b^{\mathrm{OLS}}$ & разн., \% &
$\hat\kappa_b$ & 95\% ДИ (Гарвуд) & mid-PIT\\
\midrule
2 & 51 & 50 & 0{,}5102 & 0{,}5309 & $+4{,}1$ & 1{,}921 & [1{,}426; 2{,}533] & 0{,}710\\
3 & 22 & 21 & 0{,}6148 & 0{,}6081 & $-1{,}1$ & 1{,}552 & [0{,}961; 2{,}373] & 0{,}272\\
5 & 18 & 17 & 0{,}4802 & 0{,}4666 & $-2{,}8$ & 1{,}967 & [1{,}146; 3{,}149] & 0{,}670\\
6 & 16 & 15 & 0{,}5000 & 0{,}4904 & $-1{,}9$ & 1{,}875 & [1{,}049; 3{,}092] & 0{,}594\\
7 & 9 & 8 & 0{,}7104 & 0{,}7371 & $+3{,}8$ & 1{,}251 & [0{,}540; 2{,}466] & 0{,}160\\
10 & 10 & 9 & 0{,}6612 & 0{,}6746 & $+2{,}0$ & 1{,}361 & [0{,}622; 2{,}584] & 0{,}216\\
11 & 12 & 11 & 0{,}4034 & 0{,}3617 & $-10{,}3$ & 2{,}272 & [1{,}134; 4{,}066] & 0{,}794\\
12 & 13 & 12 & 0{,}3981 & 0{,}3796 & $-4{,}6$ & 2{,}319 & [1{,}198; 4{,}050] & 0{,}820\\
13 & 12 & 11 & 0{,}4098 & 0{,}3448 & $-15{,}9$ & 2{,}237 & [1{,}117; 4{,}002] & 0{,}780\\
14 & 10 & 9 & 0{,}5025 & 0{,}5386 & $+7{,}2$ & 1{,}791 & [0{,}819; 3{,}400] & 0{,}528\\
15 & 9 & 8 & 0{,}5035 & 0{,}5054 & $+0{,}4$ & 1{,}766 & [0{,}762; 3{,}479] & 0{,}513\\
17 & 11 & 10 & 0{,}4301 & 0{,}4353 & $+1{,}2$ & 2{,}114 & [1{,}014; 3{,}887] & 0{,}717\\
18 & 7 & 6 & 0{,}6604 & 0{,}4656 & $-29{,}5$ & 1{,}298 & [0{,}476; 2{,}825] & 0{,}228\\
19 & 10 & 9 & 0{,}3161 & 0{,}2971 & $-6{,}0$ & 2{,}847 & [1{,}302; 5{,}405] & 0{,}911\\
20 & 8 & 7 & 0{,}5300 & 0{,}5713 & $+7{,}8$ & 1{,}651 & [0{,}664; 3{,}402] & 0{,}444\\
\bottomrule
\end{tabular}
\end{table}

Ниже $\emg$ лежат $11$ оснований из $15$ при концевой форме и $11$ из $15$ при
МНК --- но это \emph{разные} одиннадцать. Состав меняется: основание~18 при
концевой форме лежит выше $\emg$, а при МНК ниже, основание~20 --- наоборот.
Иными словами, даже знак отклонения не вполне устойчив к выбору формы; устойчиво
лишь их \emph{число}. Это тот же эффект, что и в \S\ref{sec:instability}, и
затрагивает он ровно те два основания, у которых выбор формы весит больше самого
отклонения.

Все пятнадцать mid-PIT-квантилей лежат внутри центральной $95\%$ области; крайние ---
$b=19$ на квантиле $0{,}911$ и $b=7$ на $0{,}160$. Основание~13, <<удачливое>>,
находится на $0{,}780$, что ничем не примечательно для одного из пятнадцати.
Все пятнадцать интервалов Гарвуда накрывают $\eg$.

\subsection{Объединённая оценка}\label{sec:pooled}

Точный условный вывод схемы~B (предложение~\ref{prop:cond}) даёт
\begin{equation}\label{eq:resB}
\hat\kappa_{\mathrm B}=\frac{M-B}{S}=1{,}8381,\qquad
95\%\ \text{ДИ}=[1{,}5939;\,2{,}1091],\qquad p=0{,}672,
\end{equation}
где интервал есть~\eqref{eq:garwood} при $K=203$, $S=110{,}4423$, а $p$-значение
получено по конвенции удвоенного минимального хвоста для
$K\sim\Pois(\eg S)$, $\eg S=196{,}71$; при mid-$p$-конвенции $p=0{,}647$.
Точное покрытие интервала~\eqref{eq:resB} равно $0{,}9540$ при $\kappa=\eg$ и не
опускается ниже $0{,}950$ ни при каком~$\kappa$.

Для сравнения, схема~A через~\eqref{eq:mle}--\eqref{eq:pivot} даёт
\begin{equation}\label{eq:resA}
\tilde\kappa=1{,}9648,\qquad 95\%\ \text{ДИ}=[1{,}7205;\,2{,}2444],\qquad p=0{,}142 .
\end{equation}
Точечная оценка~\eqref{eq:resA} на $10{,}3\%$ выше $\eg$, и именно это превышение
прежде интерпретировалось как след опущенного эвристикой фактора. При переходе к
схеме~B превышение сокращается до $3{,}2\%$, а $p$-значение растёт с $0{,}14$
до $0{,}67$.

Мы полагаем, что схема~B ближе к способу получения данных, и потому
считаем~\eqref{eq:resB} основным результатом. Оценка~\eqref{eq:resA} отвечает
остановке поиска ровно на находке и, согласно замечанию~\ref{rem:bracket}, в
смысле моментов ограничивает $\kappa$ сверху при любом правиле остановки; при
остановке <<после засухи>> истинное значение лежит ниже~\eqref{eq:resB}. Без
публикации $L_b$ проверить (A1) нельзя. Обе оценки согласуются с ЛПВ;
различаются они тем, требуют ли данные объяснения для остаточного превышения. Мы
заключаем, что не требуют. Картина в целом
показана на рис.~\ref{fig:kappa}: доверительный интервал каждого основания накрывает
и~$\eg$, и обе объединённые оценки.

\begin{figure}[t]
\centering
\includegraphics[width=.85\textwidth]{fig_kappa.pdf}
\caption{Оценки $\hat\kappa_b$ по основаниям с точными $95\%$ интервалами
Гарвуда~\eqref{eq:garwood}; площадь маркера пропорциональна $N_b$. Штриховая линия ---
значение ЛПВ $\eg$; зелёная с полосой --- объединённая оценка схемы~B и её
интервал; синяя штрихпунктирная --- оценка схемы~A $(M-1)/S$. Каждое основание по
отдельности совместимо с универсальной константой.
\encapfig{Per-base estimates $\hat\kappa_b$ with exact $95\%$ Garwood
intervals~\eqref{eq:garwood}; marker area is proportional to $N_b$. Dashed line: the
LPW value $\eg$; green with a band: the pooled scheme~B estimate and its interval;
blue dash-dotted: the scheme~A estimate $(M-1)/S$. Each base separately is
compatible with a universal constant.}}
\label{fig:kappa}
\end{figure}

\subsection{Однородность и её мощность}\label{sec:homog}

Пусть $H_0$ --- модель с единой $\kappa$, а $H_1$ --- модель со свободными $\kappa_b$.
Из условного правдоподобия предложения~\ref{prop:cond}, поскольку
$\sum_b\hat\kappa_bD_b/\ln b=K=\hat\kappa_{\mathrm B}S$, статистика
отношения правдоподобий сводится к
\begin{equation}\label{eq:LR}
\Lambda=2\sum_b K_b \ln\frac{\hat\kappa_b}{\hat\kappa_{\mathrm B}}
=2\sum_b K_b\ln\frac{K_b}{K\,s_b/S},\qquad s_b=\frac{D_b}{\ln b}.
\end{equation}
Существенно, что её нулевое распределение здесь \emph{точно}: обусловливание на
$K=\sum_bK_b$ превращает~\eqref{eq:pivotB} в мультиномиальную модель
$(K_b)\mid K\sim\mathrm{Multinomial}(K;\,s_b/S)$, и апеллировать к асимптотике
$\chi^2_{B-1}$ не нужно. Мы получаем $\Lambda=7{,}127$, точное $p=0{,}935$
($2\cdot10^6$ мультиномиальных реализаций, $\pm0{,}0002$); асимптотическое
приближение даёт $p=0{,}930$ при $14$ степенях свободы, то есть в данном случае
работает хорошо.

Однако неотвержение информативно ровно настолько, насколько мощен критерий, и здесь
следует применить к себе то же требование, что и к литературе. Мы симулировали
альтернативы, в которых $\kappa_b$ логнормально рассеяны вокруг $\eg$ с коэффициентом
вариации $cv$:

\begin{center}
\small
\begin{tabular}{lrrrrrrrrr}
\toprule
$cv$ & 0{,}05 & 0{,}10 & 0{,}15 & 0{,}20 & 0{,}25 & 0{,}30 & 0{,}35 & 0{,}40 & 0{,}50\\
мощность & 0{,}070 & 0{,}112 & 0{,}200 & 0{,}340 & 0{,}501 & 0{,}649 & 0{,}759 & 0{,}848 & 0{,}939\\
\bottomrule
\end{tabular}
\end{center}

Мощности $80\%$ критерий достигает лишь при $cv\approx0{,}37$ ($12\,000$ реализаций
на точку сетки, точность по мощности около $\pm0{,}004$). Следовательно,
$p=0{,}935$ означает: данные не требуют базозависимых частот, но и не исключают
разброса $\kappa_b$ величиной больше трети от значения константы. Утверждение
<<удачливых оснований не существует>> из имеющихся данных не следует; следует более
слабое --- <<их существование данными не поддерживается, а различия менее $\pm37\%$
принципиально неразличимы>>.

\subsection{Устойчивость}\label{sec:robust}

\paragraph{Нижнее усечение.} Мы варьировали нижнюю точку усечения, чтобы прощупать
асимптотический режим (эвристика~\eqref{eq:LPW} асимптотична, тогда как модель
постоянной интенсивности стартует от $n=2$). Результаты в таблице~\ref{tab:trunc}.
Оценка схемы~A растёт с $1{,}965$ до $2{,}099$ --- не строго монотонно (при
$n>10^{2}$ она отступает с $2{,}019$ до $2{,}012$), но систематически: рост
составляет $6{,}8\%$ при том, что все пять значений получены на вложенных
подвыборках одних и тех же данных. Это тренд, а не шум. Оценка схемы~B восходящего
тренда не показывает и остаётся вблизи $\eg$, в пределах $1{,}789$--$1{,}871$ (при
более слабом критерии выбытия $N_b\ge1$ --- в пределах $1{,}757$--$1{,}871$), а её
$p$-значения на всех уровнях усечения лежат между $0{,}53$ и $1{,}00$. Причина видна
из~\eqref{eq:pivotB}: усечение уменьшает $M$, но не $B$, тогда как схема~A
вычитает единицу независимо от объёма выборки. Мы считаем это подтверждением
диагноза~\S\ref{sec:scheme}: восходящий тренд был свойством оценки, а не данных.
Начиная с $n>10$ оценка схемы~B, напротив, монотонно \emph{убывает}
($1{,}871\to1{,}789$). Именно этого требует член второго порядка самой эвристики,
и при его учёте убывание исчезает (\S\ref{sec:second}). Обе кривые приведены на
рис.~\ref{fig:trunc}.

\begin{table}[t]
\centering
\caption{Устойчивость к нижнему усечению. Основание остаётся в выборке, если после
усечения у него сохраняется не менее \emph{двух} событий ($N_b\ge2$); при $n>10^4$ этому
критерию перестаёт удовлетворять основание~19, откуда $B=14$. Критерий существен:
при более слабом требовании $N_b\ge1$ основание~19 сохраняется, что даёт последнюю
строку таблицы. Интервал и $p_{\mathrm B}$ --- точные, схемы~B; последний столбец
($p_{\mathrm A}$, гамма-пивот схемы~A) приведён для сравнения со старой практикой.
\encaptab{Robustness to lower truncation. A base is retained if at least
\emph{two} events survive the truncation ($N_b\ge2$); at $n>10^4$ base~19 no longer
meets this criterion, whence $B=14$. The criterion matters: under the weaker
requirement $N_b\ge1$ base~19 is retained, giving the last row. The interval and
$p_{\mathrm B}$ are exact, scheme~B; the last column ($p_{\mathrm A}$, the scheme~A
Gamma pivot) is given for comparison with the old practice.}}
\label{tab:trunc}
\small
\begin{tabular}{lrrrrlrr}
\toprule
усечение & $B$ & $M$ & $(M-1)/S$ & $(M-B)/S$ & 95\% ДИ, схема B
& $p_{\mathrm B}$ & $p_{\mathrm A}$\\
\midrule
нет & 15 & 218 & 1{,}9648 & 1{,}8381 & [1{,}594; 2{,}109] & 0{,}672 & 0{,}142\\
$n>10$ & 15 & 192 & 2{,}0192 & 1{,}8712 & [1{,}606; 2{,}168] & 0{,}531 & 0{,}080\\
$n>10^2$ & 15 & 156 & 2{,}0116 & 1{,}8299 & [1{,}540; 2{,}158] & 0{,}771 & 0{,}124\\
$n>10^3$ & 15 & 125 & 2{,}0551 & 1{,}8231 & [1{,}498; 2{,}197] & 0{,}832 & 0{,}106\\
$n>10^4$ & 14 & 89 & 2{,}0990 & 1{,}7889 & [1{,}407; 2{,}242] & 1{,}000 & 0{,}116\\
\midrule
$n>10^4$, $N_b\ge1$ & 15 & 90 & 2{,}0848 & 1{,}7569 & [1{,}382; 2{,}202] & 0{,}966 & 0{,}128\\
\bottomrule
\end{tabular}
\end{table}

\begin{figure}[t]
\centering
\includegraphics[width=.75\textwidth]{fig_truncation.pdf}
\caption{Объединённая оценка при нижнем усечении: $(M-1)/S$ (схема~A) и $(M-B)/S$
(схема~B). Тренд оценки схемы~A при переходе к схеме~B исчезает.
\encapfig{The pooled estimate under lower truncation: $(M-1)/S$ (scheme~A) and
$(M-B)/S$ (scheme~B). The trend of the scheme~A estimate disappears under scheme~B.}}
\label{fig:trunc}
\end{figure}

\paragraph{Исключение по одному основанию.} Анализ с исключением одного основания даёт
$\hat\kappa_{\mathrm B}$ в пределах $[1{,}808;\,1{,}878]$ (крайние значения --- при
исключении $b=19$ и $b=3$) и точные $p$-значения в пределах $[0{,}49;\,0{,}85]$;
ни одно основание не определяет вывод.

\paragraph{Статус данных.} Ограничение основания~2 пятьюдесятью дважды проверенными
показателями переводит результат с $\tilde\kappa=1{,}965$, $p=0{,}142$ на
$\tilde\kappa=1{,}961$, $p=0{,}151$ (схема~A) и практически не меняет оценку схемы~B.
Исключение всех индексов $n>10^6$, где статус элементов --- PRP, даёт
$\tilde\kappa=1{,}951$ при $M=182$. В наихудшем сценарии, когда последний известный
PRP какого-либо основания оказывается составным, наибольший сдвиг даёт $b=17$:
$\tilde\kappa=1{,}977$. Все варианты лежат внутри интервала~\eqref{eq:resA}
и тем более внутри~\eqref{eq:resB}, поэтому возможные псевдопростые на вывод не влияют.

\section{Статистическая мощность}\label{sec:power}

Согласно~\eqref{eq:pivotB} мощность против истинной константы $\kappa=r\eg$ на уровне
$\alpha$ доступна в замкнутой форме (таблица~\ref{tab:power}). Конвенцию необходимо
оговаривать: всюду в этой работе критерий \emph{двусторонний}, с конвенцией
удвоенного минимального хвоста, $\alpha=0{,}05$. Односторонний вариант того же
критерия даёт заметно другие числа, и они приведены в правой части таблицы для
сопоставимости с работами, где конвенция не указана.

\begin{table}[t]
\centering
\caption{Мощность на уровне $\alpha=0{,}05$ для отвержения $\kappa=\eg$, когда истинная
константа превышает его в указанное число раз. Экспозиции: $S_2=26{,}02$ ($K=50$)
для одного основания~2 и $S=110{,}44$ ($K=203$) для объединения.
\encaptab{Power at level $\alpha=0.05$ for rejecting $\kappa=\eg$ when the true
constant exceeds it by the stated factor. Exposures: $S_2=26.02$ ($K=50$) for base~2
alone and $S=110.44$ ($K=203$) for the pooled sample.}}
\label{tab:power}
\small
\begin{tabular}{lrr@{\qquad}rr}
\toprule
& \multicolumn{2}{c}{двусторонний} & \multicolumn{2}{c}{односторонний}\\
\cmidrule(r){2-3}\cmidrule(l){4-5}
$\kappa/\eg$ & только Мерсенн & объединение & только Мерсенн & объединение\\
\midrule
1{,}05 & 0{,}060 & 0{,}099 & 0{,}082 & 0{,}165\\
1{,}10 & 0{,}099 & 0{,}266 & 0{,}147 & 0{,}386\\
1{,}15 & 0{,}164 & 0{,}514 & 0{,}235 & 0{,}644\\
1{,}20 & 0{,}253 & 0{,}752 & 0{,}343 & 0{,}844\\
1{,}30 & 0{,}479 & 0{,}972 & 0{,}581 & 0{,}988\\
1{,}50 & 0{,}861 & 1{,}000 & 0{,}910 & 1{,}000\\
\bottomrule
\end{tabular}
\end{table}

При мощности $80\%$ одна последовательность Мерсенна обнаруживает только отклонения
в $45{,}5\%$ и более; объединение пятнадцати оснований снижает этот порог до $21{,}3\%$
(односторонний вариант --- до $18{,}7\%$).
Достижение мощности $80\%$ против отклонения в $10\%$ требует экспозиции
$S\approx462$, то есть $K\approx824$ условных событий ($M\approx839$) при
$\kappa=\eg$ (при истинной константе, на $10\%$ большей, событий будет $\approx904$);
против $5\%$ --- $S\approx1800$ и $K\approx3210$. Функция мощности дискретного
критерия немонотонна по $S$ (пилообразна), поэтому эти пороги приведены с
точностью около $1\%$.

Эти рубежи вычислительно недостижимы, и стоит сказать насколько. Поскольку $K$ растёт
логарифмически по фронту поиска, удвоение фронтов \emph{по всем пятнадцати основаниям}
даёт $\eg\ln2\sum_b1/\ln b=9{,}64$ дополнительных события на всю объединённую выборку.
Чтобы набрать недостающие $621$ событие, нужно $621/9{,}64\approx64{,}4$ удвоения, то
есть рост фронтов в $2^{64{,}4}\approx2{,}4\cdot10^{19}$ раз; для чувствительности
в $5\%$ --- около $312$ удвоений, то есть множитель порядка $10^{94}$. Речь идёт
не о десятилетиях, а о величинах, лежащих вне досягаемости любых мыслимых вычислений.
Обе кривые мощности как функции истинного отклонения показаны на
рис.~\ref{fig:power}.

\begin{figure}[t]
\centering
\includegraphics[width=.72\textwidth]{fig_power.pdf}
\caption{Мощность двустороннего отвержения $\kappa=\eg$ на уровне $\alpha=0{,}05$.
Горизонтальная линия отмечает мощность $80\%$.
\encapfig{Power of the two-sided rejection of $\kappa=\eg$ at level $\alpha=0.05$.
The horizontal line marks $80\%$ power.}}
\label{fig:power}
\end{figure}

Отсюда правило чтения для этой литературы. Сообщение о том, что данные согласуются с ЛПВ
с точностью до нескольких процентов, не есть свидетельство в пользу ЛПВ на уровне
нескольких процентов; оно совместимо с константой, ошибочной в полтора раза.
Объединённый анализ устанавливает, однако, настоящую границу: универсальная константа,
если она существует, лежит примерно в $[1{,}59;\,2{,}11]$, и ни одно основание от неё
различимо не отклоняется.

\section{Та же ошибка в ином обличье}\label{sec:sameerror}

Счётная функция Мерсенна провоцирует родственную ошибку, которую мы приводим как
разобранный пример. Регрессия $\pi_2(t)$ на~$t$ в точках событий по $50$ дважды
проверенным показателям (нумерация $0,\dots,N-1$: наименьший показатель есть начало
отсчёта, а не событие) даёт наклон $2{,}6391$ и свободный член $-2{,}9627$ при
стандартной ошибке МНК $0{,}4469$, на вид значимый и наводящий на мысль о вторичном
члене в~\eqref{eq:LPW}. Ничего подобного.

Складываются два дефекта. Остатки кумулятивного счёта сильно автокоррелированы,
здесь $\rho_1=0{,}825$; ковариационная оценка Ньюи --- Уэста с восемью лагами
увеличивает стандартную ошибку до $0{,}7315$. Существеннее то, что регрессия использует
только точки событий и тем самым отбрасывает пустой интервал между последним показателем
и фронтом поиска. Корректная нулевая модель удерживает интервал $[t_0,T]$ фиксированным,
где $T=\ln X_{\mathrm{ver}}$, и обусловливается ровно $N-1=49$ событиями строго правее
$t_0$, что по свойству порядковых статистик пуассоновского процесса при
заданном числе событий (ср.\ доказательство предложения~\ref{prop:unbiased}(b))
означает $49$ независимых равномерных точек на $(t_0,T]$. По $60\,000$ таким реализациям свободный член имеет медиану $-1{,}518$ и
$5$--$95\%$ диапазон $[-5{,}97;\,2{,}19]$, что даёт двусторонний $p=0{,}567$ для
наблюдённого значения. Вторичный член есть артефакт обусловливания.

\subsection{Три следствия одного и того же обусловливания}

\paragraph{(i) Бутстрэп не спасает.} Естественная реакция на подозрительно малую
стандартную ошибку МНК --- бутстрэп по точкам событий. На данных основания~2
бутстрэп по $20\,000$ ресемплам даёт медиану $-2{,}949$, стандартное отклонение
$0{,}392$ и $95\%$ интервал $[-3{,}828;\,-2{,}293]$, целиком лежащий ниже нуля; доля
отрицательных значений равна $100\%$. Это не подтверждение, а повторение той же ошибки:
ресемплинг наблюдённых точек воспроизводит изменчивость внутри реализованной
конфигурации и не может восстановить изменчивость самого процесса, породившего эту
конфигурацию. Корректное нулевое распределение имеет стандартное отклонение $2{,}490$,
то есть в $6{,}4$ раза шире бутстрэпного \emph{по стандартному отклонению}
(по ширине $95\%$-интервала --- в $5{,}3$ раза).

\paragraph{(ii) Мнимая недодисперсия.} Сравнение наблюдённого $\max_i|r_i|=6{,}347$,
где $r_i=i-\hat\lambda(t_i-t_0)$, с симуляцией \emph{необусловленного} пуассоновского
процесса наводит на вывод, что счётная функция <<глаже пуассоновской>>. Против
корректного нулевого распределения --- с обусловливанием на $N-1$ событий и оценкой
$\lambda$ по последнему из них --- медиана равна $5{,}134$, а двусторонний уровень
значимости составляет $0{,}539$. Наблюдение совершенно рядовое.

\paragraph{(iii) <<Серединный дефицит>>.} Оценивание $\lambda$ по последнему событию
обращает остаток $r$ в нуль на обоих концах ($r(t_0)=r(t_N)=0$ при правильной
нумерации), превращая его траекторию в аналог броуновского моста, для которого
отрицательная экскурсия в середине диапазона возникает сама собой. Наблюдённая средняя
экскурсия на средней половине диапазона равна $-3{,}13$ при медиане симуляций
$-0{,}001$ и $p=0{,}280$; доля реализаций с отрицательной средней экскурсией
составляет $0{,}500$, то есть систематического дефицита нет. Сама траектория
остаточного процесса вместе с полосой симуляций приведена на
рис.~\ref{fig:bridge}.

\begin{figure}[t]
\centering
\includegraphics[width=.78\textwidth]{fig_resid_bridge.pdf}
\caption{Остаточный процесс $r(t)$ для основания $b=2$ вместе с полосой $5$--$95\%$
квантилей симуляций и медианой симуляций. Траектория лежит внутри полосы.
\encapfig{The residual process $r(t)$ for base $b=2$ together with the
$5$--$95\%$ quantile band of the simulations and the simulation median. The
trajectory stays inside the band.}}
\label{fig:bridge}
\end{figure}

\subsection{Основание 2 на известном фронте}

Основание~2 --- единственное, где фронт опубликован, и потому единственное, где
возможен анализ по схеме~B \emph{с использованием полной экспозиции}, а не только
условной. С $T=\ln X_{\mathrm{ver}}$ получаем
$\hat\lambda=N/(T-t_0)=2{,}7960$ против $\eg/\ln2=2{,}5695$, то есть
$\hat\kappa_2=1{,}9381$ при точном $95\%$ интервале Гарвуда $[1{,}434;\,2{,}562]$ и
пуассоновском $p=0{,}593$ (удвоенный хвост; при mid-$p$-конвенции $0{,}545$ ---
конвенцию здесь необходимо оговаривать, поскольку различие ощутимо).
Условная оценка \emph{того же среза данных} --- те же $50$ дважды проверенных
показателей --- даёт $1{,}9046$; различие невелико, поскольку последний показатель
этого среза $77\,232\,917$ случайно оказался очень близко к фронту $81\,648\,221$.
(Для сравнения, условная оценка по всем $52$ известным показателям, приведённая в
таблице~\ref{tab:bases}, равна $1{,}9215$; сопоставлять её с $\hat\kappa_2$ нельзя ---
это другая выборка.) Это совпадение одной реализации, а не свойство метода:
поправка~\S\ref{sec:cost} относится к среднему по положению последнего события
относительно фронта.

Отметим попутно, что описанный здесь режим --- ровно тот, к которому мы призываем
перейти для всех оснований: при известном $W_b$ условное обусловливание на $D_b$
становится ненужным, и точность растёт на те самые $3{,}7\%$ из \S\ref{sec:exact}.

\subsection{Теоретический вторичный член имеет противоположный знак}

Существует независимый довод против содержательной трактовки свободного члена $C_1$.
Если вторичный член в~\eqref{eq:LPW} происходит из постоянной Мертенса
$\mathfrak M=0{,}26150\ldots$, то естественная его величина есть
$(\eg/\ln2)\mathfrak M=+0{,}672$: знак положительный, тогда как наблюдается
$-2{,}963$. Ближайшие теоретические кандидаты не просто не воспроизводят наблюдаемую
величину, а предсказывают противоположный знак. Это независимо от предыдущего
подтверждает, что $C_1$ --- параметр аппроксимации, а не структуры.

\section{Структурные поправки}\label{sec:struct}

Рассмотрим более общую параметризацию интенсивности $\lambda_b=(\kappa/\ln b)\,C_b$
и четыре модели множителя:
\begin{align*}
\text{A:}&\quad C_b=1;\\
\text{B:}&\quad C_b=\prod_{p\mid b-1,\ p\ne2}\Bigl(\tfrac{p}{p-1}\Bigr)^{\alpha}
                 \prod_{p\mid b+1,\ p\ne2}\Bigl(\tfrac{p}{p-2}\Bigr)^{\beta};\\
\text{C:}&\quad C_b=\prod_{p\mid b^2-1}\Bigl(\tfrac{p}{p-1}\Bigr)^{\delta};\\
\text{D:}&\quad C_b=\exp\bigl(\theta_1\omega(b-1)+\theta_2\omega(b+1)+\theta_3\ln b\bigr).
\end{align*}
Показатель в модели~C обозначен через $\delta$, а не $\gamma$, чтобы не смешивать его
с постоянной Эйлера. Правдоподобие здесь --- \emph{условное пуассоновское}
предложения~\ref{prop:cond},
$\ell_{\mathrm B}=\sum_b\bigl[K_b\ln\mu_b-\mu_b-\ln K_b!\bigr]$ с
$\mu_b=\kappa C_b D_b/\ln b$; в предыдущей редакции использовалось
правдоподобие~\eqref{eq:loglik} схемы~A, что несовместимо с
\S\ref{sec:scheme}. Результаты в таблице~\ref{tab:struct}. Ни одна параметризация не
даёт значимого улучшения; байесовские факторы по BIC при $n=K=203$ подавляюще
поддерживают базовую модель.

\begin{table}[t]
\centering
\caption{Сравнение структурных моделей под условным пуассоновским правдоподобием
предложения~\ref{prop:cond} (полная форма, с $-\ln K_b!$); LR и $p$ даны относительно
модели~A; BIC вычислен при $n=K=203$.
\encaptab{Comparison of structural models under the conditional Poisson likelihood of
Proposition~\ref{prop:cond} (full form, including $-\ln K_b!$); LR and $p$ are
relative to model~A; BIC is computed with $n=K=203$.}}
\label{tab:struct}
\small
\begin{tabular}{lrrrrrrr}
\toprule
Модель & параметров & $\ln L$ & LR & $p$ & AIC & BIC & $P(\text{модель}\mid\text{данные})$\\
\midrule
A & 1 & $-35{,}753$ & --- & --- & 73{,}51 & 76{,}82 & 0{,}906\\
B & 3 & $-35{,}162$ & 1{,}18 & 0{,}554 & 76{,}32 & 86{,}26 & 0{,}008\\
C & 2 & $-35{,}458$ & 0{,}59 & 0{,}443 & 74{,}92 & 81{,}54 & 0{,}085\\
D & 4 & $-34{,}797$ & 1{,}91 & 0{,}591 & 77{,}59 & 90{,}85 & 0{,}001\\
\bottomrule
\end{tabular}
\end{table}

\begin{remark}[Чувствительность BIC]
Выбор эффективного объёма выборки для BIC неочевиден: правдоподобие имеет
$B=15$ независимых слагаемых, но информация растёт как $K=203$. При $n=15$
апостериорные вероятности составляют $0{,}661$ / $0{,}080$ / $0{,}229$ / $0{,}030$,
при $n=203$ --- $0{,}906$ / $0{,}008$ / $0{,}085$ / $0{,}001$. Вывод в пользу модели~A
при обоих выборах, но при $n=203$ он существенно сильнее. Мы отчитываемся по
$n=K$ как соответствующему числу независимых условных событий.
\end{remark}

\paragraph{Фиксированная арифметическая поправка.} Отдельно рассмотрена модель с
\emph{фиксированным} (не оцениваемым) множителем
$C^{\mathrm{fix}}_b=\prod_{p\mid b-1,p\ne2}\frac{p}{p-1}\prod_{p\mid b+1,p\ne2}\frac{p}{p-2}$,
среднее значение которого по пятнадцати основаниям равно $2{,}462$. Она даёт
$\hat\kappa=0{,}759$. Подчеркнём, что сравнивать это значение с $\eg$ бессмысленно:
в перепараметризованной модели $\kappa$ имеет иной масштаб, и такой <<критерий
согласия>> проверял бы лишь равенство среднего $C_b$ единице. Поскольку модели A и
$C^{\mathrm{fix}}$ ненестированы и имеют одинаковое число свободных параметров,
корректное сравнение --- прямо по правдоподобию: $\ln L=-51{,}53$ против $-35{,}75$,
разность $-15{,}78$. Фиксированная поправка решительно отвергается.

\paragraph{Логарифмические поправки.} Здесь необходимо исправление по существу.
Модель $\lambda_b(t)=(\kappa/\ln b)(1+c_1/t+c_2/t^2)$, использованная в предыдущей
редакции, \emph{не является допустимой интенсивностью}: при полученных там оценках
$\hat c_1=+0{,}108$, $\hat c_2=-0{,}826$ множитель обращается в нуль при
$t=0{,}8564$ и отрицателен левее, то есть при $n<2{,}355$; в частности, при
$t_0=\ln 2$ он равен $-0{,}5634$. Основания с $n_{\min}=2$ ($b=2,6,10,12,18$)
интегрируют интенсивность по области, где она отрицательна, и мера
$\int\lambda\,dt$ на этом участке отрицательна. То, что правдоподобие всё же
вычислялось, --- случайность: наименьший индекс, являющийся событием, равен $3$,
то есть $t=1{,}0986>0{,}8564$, так что под логарифм отрицательные значения не
попали.

Правильная параметризация того же семейства ---
\begin{equation}\label{eq:logcorr}
\lambda_b(t)=\frac{\kappa}{\ln b}\exp\Bigl(\frac{c_1}{t}+\frac{c_2}{t^2}\Bigr),
\end{equation}
положительная тождественно и вкладывающаяся в нулевую модель при $c_1=c_2=0$.
Она оценивается условным правдоподобием точечного
процесса~\eqref{eq:condinhom}: сумма $\sum\ln\lambda(t_i)$ берётся по событиям
\emph{без последнего}, а интеграл $\int\lambda$ --- от $t_0$ до последнего события
$t_0+D_b$, а не до неизвестного фронта (замечание~\ref{rem:inhom}; правдоподобие
одного лишь счёта здесь недопустимо: оно отбрасывает положения событий, которыми
только и определяются $c_1,c_2$). Результат: $\hat\kappa=1{,}748$,
$\hat c_1=0{,}75$, $\hat c_2=-1{,}44$; отношение правдоподобий против $c_1=c_2=0$
равно $2{,}03$ при двух степенях свободы, $p=0{,}36$ (тег~[S17]). Значимого улучшения
нет. Заметим, однако, что $\hat c_1$ близко к значению $\ln2$, которое
предсказывает сама эвристика; этот сюжет разобран в~\S\ref{sec:second}.

\section{Член второго порядка самой эвристики}\label{sec:second}

До сих пор проверялась модель постоянной интенсивности~\eqref{eq:lam}. Но сама
эвристика ЛПВ предсказывает интенсивность, \emph{не} постоянную при конечных~$n$,
и это меняет смысл сравнения $\hat\kappa$ с $\eg$.

\paragraph{Вывод.} Пусть $p$ простое, $p\nmid b-1$, и $q$ --- простой делитель
$\Rb(p)$. Тогда $b$ имеет порядок $p$ по модулю $q$, откуда $q\equiv1\pmod p$, а
так как $\Rb(p)$ нечётно, $q\equiv1\pmod{2p}$. Значит, у $\Rb(p)$ нет простых
делителей меньше $2p+1$, и по теореме Мертенса вероятность простоты
увеличивается в $\prod_{q<2p}(1-1/q)^{-1}\approx\eg\ln(2p)$ раз по сравнению со
случайным числом той же величины. Это и есть вывод~\cite{wagstaff1983,t5k}:
\begin{equation}\label{eq:probp}
\Pr\bigl(\Rb(p)\ \text{простое}\bigr)\approx\frac{\eg\ln(a_bp)}{p\ln b},\qquad a_b\ge2.
\end{equation}
Уточнение Вагстаффа для $b=2$ учитывает, что $q\equiv\pm1\pmod8$, и даёт $a=2$ при
$p\equiv3\pmod4$ и $a=6$ при $p\equiv1\pmod4$. Для общего $b$ аналогичное
необходимое условие --- $b$ есть квадратичный вычет по модулю~$q$, поскольку
$b^p\equiv1$ и $p$ нечётно; наименьший допустимый кандидат $2kp+1$ определяет
$a_b=2k$. Умножая~\eqref{eq:probp} на плотность простых $1/\ln n$ и переходя к
$t=\ln n$, получаем
\begin{equation}\label{eq:lam2}
\lambda_b(t)=\frac{\kappa}{\ln b}\Bigl(1+\frac{c_b}{t}\Bigr),\qquad
c_b=\langle\ln a_b\rangle,\qquad \kappa=\eg,
\end{equation}
а для счётной функции --- $\pi_b(x)\approx(\eg/\ln b)\,(\ln x+c_b\ln\ln x)$.
Асимптотика~\eqref{eq:LPW} от этого не меняется, но при $t\approx10$--$18$ поправка
$c_b/t$ составляет от $4$ до $17\%$ --- столько же или больше, чем обсуждаемые в
работе эффекты. Усреднение по простым $p\in[10^5,2\cdot10^5]$ даёт $c_b$ от $1{,}23$
($b=19$) до $1{,}74$ ($b=15$), для $b=2$ --- $1{,}242=\tfrac12(\ln2+\ln6)$, как и
должно быть (скрипт \texttt{lpw\_second\_order.py}, тег~[S14]). Значение
$c=\ln2\approx0{,}693$ --- нижняя граница, использующая лишь $q\equiv1\pmod{2p}$.
Подчеркнём, что~\eqref{eq:probp} --- эвристика; её член второго порядка менее
надёжен, чем главный, и мы рассматриваем $c$ как внешний параметр в диапазоне
$[\ln2;\ \langle\ln a_b\rangle]$.

\paragraph{Точный условный вывод при заданном $c$.} По
замечанию~\ref{rem:inhom} условно на $D_b$ внутренние события образуют
пуассоновский процесс интенсивности~\eqref{eq:lam2} на $(t_0,t_0+D_b)$. Поскольку
$\int_{t_0}^{t_0+D_b}(1+c/t)\,dt=D_b+c\ln\bigl((t_0+D_b)/t_0\bigr)$,
предложение~\ref{prop:cond} переносится дословно:
\begin{equation}\label{eq:pivotc}
K\mid(D_b)_b\sim\Pois\bigl(\kappa S(c)\bigr),\qquad
S(c)=\sum_b\frac{D_b+c_b\ln\bigl((t_0+D_b)/t_0\bigr)}{\ln b},
\end{equation}
оценка есть $K/S(c)$, и интервал Гарвуда~\eqref{eq:garwood} применим с заменой $S$
на $S(c)$.

\begin{table}[t]
\centering
\caption{Точный условный вывод при трёх значениях члена второго
порядка~\eqref{eq:lam2}: оценка $\hat\kappa=K/S(c)$ по~\eqref{eq:pivotc} и точное
двустороннее $p$ против $\kappa=\eg$. <<В>> --- $c_b=\langle\ln a_b\rangle$ по
Вагстаффу. Усечение, как в таблице~\ref{tab:trunc}. Тег~[S16].
\encaptab{Exact conditional inference for three values of the second-order
term~\eqref{eq:lam2}: the estimate $\hat\kappa=K/S(c)$ from~\eqref{eq:pivotc} and the
exact two-sided $p$ against $\kappa=\eg$. ``W'' is Wagstaff's
$c_b=\langle\ln a_b\rangle$. Truncation as in Table~\ref{tab:trunc}. Tag~[S16].}}
\label{tab:second}
\small
\begin{tabular}{lrrrrrrr}
\toprule
усечение & $K$ & $\hat\kappa\,(c{=}0)$ & $\hat\kappa\,(c{=}\ln2)$ & $\hat\kappa$\,(В)
& $p\,(c{=}0)$ & $p\,(c{=}\ln2)$ & $p$\,(В)\\
\midrule
нет & 203 & 1{,}838 & 1{,}625 & 1{,}453 & 0{,}672 & 0{,}200 & 0{,}003\\
$n>10$ & 177 & 1{,}871 & 1{,}712 & 1{,}578 & 0{,}531 & 0{,}630 & 0{,}111\\
$n>10^2$ & 141 & 1{,}830 & 1{,}704 & 1{,}595 & 0{,}771 & 0{,}636 & 0{,}201\\
$n>10^3$ & 110 & 1{,}823 & 1{,}715 & 1{,}620 & 0{,}832 & 0{,}738 & 0{,}344\\
$n>10^4$ & 75 & 1{,}789 & 1{,}695 & 1{,}612 & 1{,}000 & 0{,}723 & 0{,}421\\
\bottomrule
\end{tabular}
\end{table}

\paragraph{Результаты.} Таблица~\ref{tab:second} даёт четыре вывода.

\emph{(i) Член второго порядка данными не определяется.} Профильная ОМП общего $c$
по условному правдоподобию~\eqref{eq:condinhom} на всех данных равна $-0{,}40$ с
$95\%$ интервалом $[-0{,}69;\,0{,}44]$ (левый конец --- граница допустимости
$1+c/t_0>0$). Но уже при $n>10$ интервал
расширяется до $[-1{,}68;\,2{,}48]$, а при $n>10^3$ --- до $[-5{,}3;\,20]$
(тег~[S16]). Положения событий внутри последовательностей почти не несут
информации о $c$: сдвиг $c$ компенсируется сдвигом $\kappa$. Двухпараметрическая
модель~\eqref{eq:logcorr} также не даёт значимого улучшения (\S\ref{sec:struct}).
Значение $c$ поэтому должно браться из теории, а не из данных.

\emph{(ii) Вывод о $\kappa$ от него зависит сильнее, чем от схемы наблюдения.}
При $n>10$ переход от $c=0$ к $c=\ln2$ снижает оценку на $8{,}5\%$, к
$c=\langle\ln a_b\rangle$ --- на $15{,}7\%$. Это больше цены выбора схемы
наблюдения ($6{,}9\%$, \S\ref{sec:cost}) и вдвое-втрое больше остаточного разрыва
$(M-B)/S$ с $\eg$ ($3{,}2\%$). Сравнение оценки модели постоянной интенсивности с
$\eg$ --- не проверка ЛПВ при конечных $n$: сама ЛПВ предсказывает, что такая
оценка \emph{выше} $\eg$.

\emph{(iii) ЛПВ со своим членом второго порядка согласуется с данными.} При
$c=\ln2$ и $n>10$ точное $p$ лежит в пределах $0{,}63$--$0{,}74$. При
$c=\langle\ln a_b\rangle$ оно равно $0{,}11$--$0{,}42$. На полной выборке с
$c=\langle\ln a_b\rangle$ получается $p=0{,}003$. Но полная выборка включает
$n\le10$, где разложение~\eqref{eq:lam2} бессмысленно: при $n=2$ множитель
$1+c/t$ равен $2$ даже при $c=\ln2$, а сами индексы $2,3,5,7$ продиктованы
арифметикой малых чисел. Этот результат мы не считаем свидетельством против ЛПВ;
он предостерегает от использования малых индексов при $c\ne0$.

\emph{(iv) Тренд при усечении объясняется членом второго порядка.} При $c=0$
оценка монотонно убывает с ростом нижней границы, с $1{,}871$ при $n>10$ до
$1{,}789$ при $n>10^4$ ($-4{,}4\%$). Именно этого ожидают от~\eqref{eq:lam2}:
среднее $1/t$ по экспозиции падает с ростом $t_0$. При $c=\ln2$ те же четыре
оценки лежат в пределах $1{,}695$--$1{,}715$ (размах $1{,}2\%$). Подвыборки
вложены, поэтому это наблюдение иллюстративно, а не является отдельным критерием.

\paragraph{Следствие для фольклора.} Член $c_b\ln\ln x$ в счётной функции объясняет
и наблюдение, с которого начиналась работа: нормированные частоты лежат ниже
$\emg$ и медленно сходятся к нему снизу. При~\eqref{eq:lam2} наклон $t_k$ по $k$
равен $\emg\ln b/(1+c_b\langle1/t\rangle)$, так что нормированная частота
занижена в $1+c_b\langle1/t\rangle$ раз и стремится к $\emg$ как $1/t$, то есть
как $1/\ln n$ --- действительно медленно. В~\S\ref{sec:notbias} показано, что
подгонка к правому концу эффекта не создаёт; здесь видно, откуда он берётся:
из самой эвристики, а не из~$\kappa\ne\eg$.

\paragraph{Итог.} Вопрос <<совместимы ли данные с ЛПВ>> распадается на два:
о главном члене, $\kappa=\eg$, и о члене второго порядка, $c_b$. Второй
данными не определяется, но влияет на ответ на первый сильнее, чем всё
разобранное в \S\ref{sec:scheme}. При любом $c\in[\ln2;\,\langle\ln a_b\rangle]$
и $n>10$ ЛПВ не отвергается. Уточнение члена второго порядка --- задача теории
чисел, а не статистики. Без него количественная проверка ЛПВ точнее
$\pm10\%$ недостижима и при росте фронтов: член $c_b/t$ убывает лишь как
$1/\ln n$.

\section{Заключение}

\begin{enumerate}
\item Употребительные оценки нормированной частоты обобщённых репьюнитных простых
\emph{не} смещены схемой подгонки: и концевая форма, и наклон МНК несмещены точно при
схеме~A и с экспоненциально малой по $\mu$ погрешностью при схеме~B (не более
$0{,}73\%$ для концевой формы и $2{,}3\%$ для МНК при $\mu\ge7$).
Смещена лишь обратная величина, и это следствие выпуклости. Расхожее объяснение
<<медленной сходимости снизу>> через цензурирование несостоятельно при обеих схемах.

\item Настоящий дефект употребительной практики --- в том, что различные формы оценки
расходятся до $29{,}5\%$, а точечные значения публикуются без доверительных интервалов
при том, что стандартные ошибки МНК в этой задаче занижены: поправка Ньюи ---
Уэста увеличивает их в $1{,}6$ раза, а корректное условное нулевое распределение
шире номинального в $5{,}6$ раза (\S\ref{sec:sameerror}).
Классификация оснований на <<удачливые>> и <<неудачливые>> переворачивается от смены
формы оценки на неизменных данных; у двух оснований из пятнадцати переворачивается
даже знак отклонения от $\emg$.

\item Главный результат --- точное условное правдоподобие схемы~B. Незаписанные
фронты суть мешающие параметры Неймана --- Скотта, и профильная оценка $M/S$
несостоятельна. Положение последнего события $D_b$ достаточно для фронта, и
условно на нём величина $N_b-1$ распределена как $\Pois(\lambda_bD_b)$ при любом
не зависящем от находок правиле выбора фронта. Тот же результат даёт
интегрирование по фронту с инвариантной мерой $dW_b$. Отсюда $M-B\mid D\sim\Pois(\kappa S)$:
правильным пивотом при фиксированном фронте является пуассоновский, а не
$\Gam(M,1)$, правильной объединённой оценкой --- $(M-B)/S$, а не $(M-1)/S$, и
доверительные интервалы получаются в замкнутой форме без симуляционной инверсии.
Условное правдоподобие стоит одного события на каждое основание, а не одного
события в сумме, --- точный аналог потери одной степени свободы на группу в
примере Неймана --- Скотта. Различение таких схем восходит к теории надёжности; новым здесь
является случай недокументированного фронта, когда экспозицию нельзя взять из
протокола и её приходится восстанавливать из находок
(замечание~\ref{rem:priority}).

\item Это даёт $\kappa=1{,}838$ при точном $95\%$ интервале $[1{,}594;\,2{,}109]$ и
$p=0{,}672$ против $\eg$. Сообщавшееся десятипроцентное превышение над ЛПВ, включая
наше собственное в более ранней редакции этой работы, есть артефакт схемы наблюдения.
Старый гамма-пивот при схеме~B отвергает верную гипотезу в $11$--$17\%$ случаев.
Переход к схеме~B независимо устраняет и систематический тренд объединённой оценки
при усечении по $n_{\min}$.

\item Вывод держится на одном допущении (A1): фронт выбирался без оглядки на
находки. Проверить его данными нельзя: $(M-1)/S$ и $(M-B)/S$ суть функции одних и
тех же статистик и различаются на $6{,}9\%$. Без (A1) строго сохраняется лишь
верхняя граница (тождество Вальда); правило <<остановиться после долгой полосы
без находок>> смещает вверх и $(M-B)/S$, так что двусторонней вилки нет.
Ни одно из разобранных нарушений (A1) не создаёт превышения над ЛПВ. Четыре
документированных фронта указывают, если на что-то, на отклонение в сторону
остановки вблизи находки, а не <<после засухи>>; гибридная оценка с известными
фронтами $b=2$ и $b=10$ равна $1{,}853$.
Практическое следствие --- необходимость публиковать фронты поиска
$L_b$ (а вместе с ними и протоколы полноты просеивания), как это делает GIMPS для
$b=2$. До тех пор любой объединённый анализ этих последовательностей опирается
на непроверяемое допущение, цена которого --- не менее $7\%$, то есть вдвое больше
остаточного разрыва с ЛПВ.

\item Гипотеза ЛПВ имеющимися данными не отвергается, а базозависимые поправки не
требуются. Но и то и другое --- утверждения о низкой чувствительности критериев:
условный мультиномиальный критерий однородности различает лишь разброс $\kappa_b$
порядка $\pm37\%$, а объединённый критерий --- отклонения константы не менее
$21{,}3\%$ (двусторонний вариант). Достижение чувствительности в $10\%$ потребовало
бы роста фронтов поиска примерно в $2\cdot10^{19}$ раз.

\item Сама эвристика ЛПВ предсказывает при конечных $n$ интенсивность
$(\kappa/\ln b)(1+c_b/t)$ с $c_b\in[\ln2;\,1{,}74]$, и сравнение оценки модели
постоянной интенсивности с $\eg$ не является проверкой ЛПВ. Член второго порядка
данными не определяется, но сдвигает оценку $\kappa$ на $8$--$16\%$; при его
учёте ЛПВ согласуется с данными, тренд при усечении исчезает, а <<медленная
сходимость снизу>> получает объяснение. Уточнение $c_b$ --- задача теории чисел,
без решения которой проверка ЛПВ точнее $\pm10\%$ недостижима.
\end{enumerate}


\appendix
\section{Верификация исходных данных}\label{app:data}

Все двадцать последовательностей (пятнадцать основного набора и пять из проверки
границы отбора) сверены с записями OEIS 27 августа 2026~г.
Поскольку в объединённый анализ входят только тройки $(N_b,n_{\min},n_{\max})$,
в таблице~\ref{tab:verify} для основного набора приведены число известных простых, из которого
$N_b$ получается вычитанием единицы, и оба крайних индекса; полные списки
индексов содержатся в скрипте \texttt{analysis\_v2.py}.

\begin{table}[h]
\centering
\caption{Сверка с OEIS. Столбец <<последнее пополнение>> указывает происхождение
наибольшего известного индекса.
\encaptab{Cross-check against OEIS. The column ``last addition'' gives the
provenance of the largest known index.}}
\label{tab:verify}
\small
\begin{tabular}{rlrrrl}
\toprule
$b$ & OEIS & простых & $n_{\min}$ & $n_{\max}$ & последнее пополнение\\
\midrule
2 & A000043 & 52 & 2 & 136\,279\,841 & GIMPS, окт.\ 2024\\
3 & A028491 & 23 & 3 & 8\,530\,117 & P.~Bourdelais\\
5 & A004061 & 19 & 3 & 3\,300\,593 & P.~Bourdelais\\
6 & A004062 & 17 & 2 & 3\,360\,347 & P.~Bourdelais\\
7 & A004063 & 10 & 5 & 1\,264\,699 & P.~Bourdelais\\
10 & A004023 & 11 & 2 & 8\,177\,207 & Propper и Batalov, май 2021\\
11 & A005808 & 13 & 17 & 1\,868\,983 & P.~Bourdelais\\
12 & A004064 & 14 & 2 & 769\,543 & P.~Bourdelais, дек.\ 2014\\
13 & A016054 & 13 & 5 & 1\,503\,503 & P.~Bourdelais\\
14 & A006032 & 11 & 3 & 1\,724\,417 & P.~Bourdelais, 13.07.2026\\
15 & A006033 & 10 & 3 & 639\,833 & P.~Bourdelais, 22.04.2019\\
17 & A006034 & 12 & 3 & 1\,990\,523 & P.~Bourdelais, 03.08.2020\\
18 & A133857 & 8 & 2 & 1\,270\,141 & P.~Bourdelais\\
19 & A006035 & 11 & 19 & 209\,359 & P.~Bourdelais, 27.08.2010\\
20 & A127995 & 9 & 3 & 984\,349 & P.~Bourdelais\\
\bottomrule
\end{tabular}
\end{table}

Отдельно предупредим читателя, сверяющего данные по вторичным источникам: сводные
таблицы (страница <<Repunit primes>> в OEIS Wiki, статья <<List of repunit primes>>
в Википедии) обновляются с запаздыванием и на момент написания работы отставали
от самих записей OEIS на один член у оснований $b=14,15,17,18,20$. Авторитетным
источником следует считать записи OEIS.

\subsection*{Независимый пересчёт префиксов}

Сверка с OEIS подтверждает, что перечисленные индексы дают простое, но не
гарантирует \emph{полноты}: пропущенный член сдвинул бы $N_b$, а с ним и обе
объединённые оценки. Этот пробел мы закрыли прямым пересчётом. Для каждого из
двадцати оснований $2\le b\le26$, не являющихся точными степенями, независимой
реализацией\footnote{Гибридный CPU/GPU поисковик: сегментное сито по малым
делителям, trial factoring в арифметике Монтгомери, PRP через GMP и GWNUM/IBDWT.
Исходный код и протокол прогонов --- в том же репозитории, скрипты
\texttt{verify\_sequences.sh} и \texttt{compare\_verify.py}.} проверены
\emph{все} индексы $n<10^4$ подряд, а не только перечисленные в OEIS. Каждый
отрицательный вердикт пересчитан вторым, независимым бэкендом (сплошной
double-check, как в GIMPS), чтобы исключить потерянную находку.

Результат: $153$ члена подтверждены, лишних нет, пропущенных нет, расхождений
между бэкендами нет --- по всем двадцати основаниям. Покрытие по основаниям
неравномерно (от одного члена при $b=18$ до двадцати двух при $b=2$): на
основной набор $b\le20$ приходится $128$ подтверждённых индексов из $233$
известных. Выше $10^4$ полный пересчёт вычислительно недоступен, и там мы
по-прежнему опираемся на OEIS.

Отметим, что этот пересчёт не был формальностью: в первом прогоне независимая
реализация объявила $R_{3181}(23)$ составным, разойдясь с A204940. Разбор показал
ошибку в самой реализации --- всплеск ошибки округления FFT на последней итерации,
после последней периодической проверки, --- а не в данных OEIS. После исправления
расхождений не осталось. Мы приводим этот эпизод как иллюстрацию к тезису
\S\ref{sec:scheme}: ложное <<составное>> невидимо без повторного счёта, и именно
поэтому фронты и протоколы проверки должны публиковаться.

\section*{Доступность данных и кода}
Все использованные последовательности приведены в таблице~\ref{tab:data}.

\emph{Данные.} Скрипт \texttt{fetch\_data.sh} загружает b-файлы OEIS для всех
пятнадцати последовательностей основного набора, \texttt{fetch\_extra.sh} --- для
пяти оснований $21\le b\le26$, использованных в проверке границы отбора
(\S\ref{sec:bases}). Загруженные файлы включены в репозиторий, так что анализ
воспроизводится без обращения к сети.

\emph{Анализ.} Скрипт \texttt{analysis\_v2.py} (numpy, scipy, matplotlib); каждое
число настоящей работы помечено в его выводе (\texttt{results\_v2.txt}) тегом
вида~[S3], соответствующим разделу или таблице. Объединённый анализ, все
доверительные интервалы, критерий однородности, расчёты мощности и структурные
модели требуют только троек $(N_b,n_{\min},n_{\max})$, зашитых в скрипт, и
воспроизводятся без обращения к данным; полные списки индексов нужны лишь для
столбца $\hat G^{\mathrm{OLS}}_b$ таблицы~\ref{tab:bases}, для \S\ref{sec:sameerror}
и для модели~\eqref{eq:logcorr}. Генератор случайных чисел инициализирован
фиксированным зерном ($20260914$), поэтому все величины Монте-Карло воспроизводятся
побитово. Скрипт также строит все четыре рисунка. Таблицу~\ref{tab:stop}
(правила остановки, нарушающие (A1)) строит отдельный скрипт
\texttt{stopping\_rules.py} (тег~[S13], зерно $20260924$). Член второго
порядка (\S\ref{sec:second}), модель~\eqref{eq:logcorr} и критерий однородности при
$B=20$ вычисляет \texttt{lpw\_second\_order.py} (теги~[S14]--[S17]), проверку (A1)
по документированным фронтам и гибридную оценку --- \texttt{frontiers.py}
(тег~[S18], источники фронтов перечислены в самом скрипте). Вывод всех трёх
скриптов сохранён в \texttt{results\_v3.txt}.
% ВНИМАНИЕ: analysis_v2.py и results_v2.txt (версия 2) в репозитории
% отсутствуют; их нужно выложить, иначе ссылки на теги [S0]-[S12] висячие.

\emph{Независимый пересчёт последовательностей.} \texttt{verify\_sequences.sh}
запускает поисковик по всем индексам $n<10^4$ для каждого основания,
\texttt{compare\_verify.py} сверяет результат с b-файлами OEIS и возвращает
ненулевой код при любом расхождении (приложение~\ref{app:data}).

\emph{Оговорка о величинах Монте-Карло.} Основной результат~\eqref{eq:resB},
интервалы таблицы~\ref{tab:bases}, таблица~\ref{tab:trunc} и
таблица~\ref{tab:power} \emph{детерминированы}: они получены точной инверсией
пуассоновского распределения и от генератора случайных чисел не зависят. От
симуляции зависят лишь таблицы~\ref{tab:mc}, \ref{tab:calib} и~\ref{tab:stop}, точное
$p$-значение критерия однородности и кривая его мощности, а также нулевые
распределения \S\ref{sec:sameerror}; их точность указана на месте. Это отличие от
предыдущей редакции существенно: там симуляционной была и сама основная оценка
интервала.

Репозиторий: \url{https://github.com/Wilps93/repunit-hunt} --- каталог
\texttt{paper/}. Дерево файлов и порядок запуска описаны в \texttt{paper/README.md}.

% п. 4.14 ГОСТ Р 7.0.7-2021: перечень затекстовых библиографических ссылок
% предваряют словами «Список источников»; babel-russian по умолчанию печатает
% «Список литературы», поэтому \refname переопределяем.
\renewcommand{\refname}{Список источников}
\begin{thebibliography}{99}

\bibitem{wagstaff1983} Wagstaff S.~S., Jr. Divisors of Mersenne numbers~//
Mathematics of Computation. 1983. Vol.~40, №~161. P.~385--397.
DOI: 10.1090/S0025-5718-1983-0679454-X.

\bibitem{pomerance1981} Pomerance C. Recent developments in primality testing~//
The Mathematical Intelligencer. 1981. Vol.~3, №~3. P.~97--105.
DOI: 10.1007/BF03022861.

\bibitem{t5k} Heuristics: deriving the Wagstaff Mersenne conjecture
[Электронный ресурс] // The PrimePages: prime number research and records.
URL: \url{https://t5k.org/mersenne/heuristic.html} (дата обращения: 27.08.2026).

\bibitem{dubner1993} Dubner H. Generalized repunit primes~//
Mathematics of Computation. 1993. Vol.~61, №~204. P.~927--930.
DOI: 10.1090/S0025-5718-1993-1185243-9.

\bibitem{bourdelais2009} Bourdelais P. A generalized repunit conjecture
[Электронный ресурс]~// NMBRTHRY mailing list. 25 June 2009.
URL: \url{https://listserv.nodak.edu/cgi-bin/wa.exe?A2=NMBRTHRY;417ab0d6.0906}
(дата обращения: 25.08.2026).

\bibitem{oeis} The On-Line Encyclopedia of Integer Sequences
[Электронный ресурс] / OEIS Foundation Inc.
URL: \url{https://oeis.org} (дата обращения: 27.08.2026).

\bibitem{wagstaff2025} Wagstaff S.~S., Jr. Two Mersenne prime conjectures
[Электронный ресурс]~// Journal of Integer Sequences. 2025. Vol.~28.
Article 25.7.2.
URL: \url{https://cs.uwaterloo.ca/journals/JIS/VOL28/Wagstaff/wag14.html}
(дата обращения: 24.09.2026).

\bibitem{dominguez2026} Dominguez J. Arithmetic bias in the distribution of Mersenne
prime exponents and the divisor structure of $p-1$ [Электронный ресурс] // arXiv.
2026. arXiv:2603.08994. URL: \url{https://arxiv.org/abs/2603.08994}
(дата обращения: 27.08.2026).

\bibitem{neymanscott1948} Neyman J., Scott E.~L. Consistent estimates based on
partially consistent observations~// Econometrica. 1948. Vol.~16, №~1. P.~1--32.
DOI: 10.2307/1914288.

\bibitem{andersen1970} Andersen E.~B. Asymptotic properties of conditional
maximum-likelihood estimators~// Journal of the Royal Statistical Society.
Series~B. 1970. Vol.~32, №~2. P.~283--301.
DOI: 10.1111/j.2517-6161.1970.tb00842.x.

\bibitem{jha2026} Jha A. The Poisson tail conjecture for primes in short intervals
[Электронный ресурс] // arXiv. 2026. arXiv:2605.23014.
URL: \url{https://arxiv.org/abs/2605.23014} (дата обращения: 27.08.2026).

\bibitem{gimps} PrimeNet [Электронный ресурс] / Great Internet Mersenne Prime
Search (GIMPS). Программное обеспечение Prime95/MPrime, версия~30.8.
URL: \url{https://www.mersenne.org/} (дата обращения: 27.08.2026;
отчёт \emph{Milestones} считан 27.08.2026 04:30~UTC).

\bibitem{epstein1954trunc} Epstein B. Truncated life tests in the exponential
case~// The Annals of Mathematical Statistics. 1954. Vol.~25, №~3. P.~555--564.
DOI: 10.1214/aoms/1177728723.

\bibitem{epsteinsobel1954} Epstein B., Sobel M. Some theorems relevant to life
testing from an exponential distribution~// The Annals of Mathematical
Statistics. 1954. Vol.~25, №~2. P.~373--381. DOI: 10.1214/aoms/1177728793.

\bibitem{bartholomew1963} Bartholomew D.~J. The sampling distribution of an
estimate arising in life testing~// Technometrics. 1963. Vol.~5, №~3.
P.~361--374. DOI: 10.1080/00401706.1963.10490104.

\bibitem{sengupta2009} Sengupta S.~N., Sinha B.~K., Bagchi S.~B. Estimation of
mean life and reliability for exponential life distribution using time censored
sample data~// Statistics \& Probability Letters. 2009. Vol.~79, №~7.
P.~899--906. DOI: 10.1016/j.spl.2008.11.026.

\bibitem{kivela2015} Kivel\"{a} M., Porter M.~A. Estimating interevent time
distributions from finite observation periods in communication networks~//
Physical Review~E. 2015. Vol.~92, №~5. Article~052813.
DOI: 10.1103/PhysRevE.92.052813.

\bibitem{straussSadler1989} Strauss D., Sadler P.~M. Classical confidence
intervals and Bayesian probability estimates for ends of local taxon ranges~//
Mathematical Geology. 1989. Vol.~21, №~4. P.~411--427. DOI: 10.1007/BF00897326.

\bibitem{marshall1990} Marshall C.~R. Confidence intervals on stratigraphic
ranges~// Paleobiology. 1990. Vol.~16, №~1. P.~1--10.
DOI: 10.1017/S0094837300009672.

\bibitem{solow1993} Solow A.~R. Inferring extinction from sighting data~//
Ecology. 1993. Vol.~74, №~3. P.~962--964. DOI: 10.2307/1940821.

\bibitem{robsonwhitlock1964} Robson D.~S., Whitlock J.~H. Estimation of a
truncation point~// Biometrika. 1964. Vol.~51, №~1--2. P.~33--39.
DOI: 10.1093/biomet/51.1-2.33.

\bibitem{liu2023} Liu Z., Bhandaram U., Garg N. Quantifying spatial
under-reporting disparities in resident crowdsourcing~// Nature Computational
Science. 2024. Vol.~4, №~1. P.~57--65. DOI: 10.1038/s43588-023-00572-6.

\bibitem{lancaster2000} Lancaster T. The incidental parameter problem since
1948~// Journal of Econometrics. 2000. Vol.~95, №~2. P.~391--413.
DOI: 10.1016/S0304-4076(99)00044-5.

\bibitem{kalbfleischsprott1970} Kalbfleisch J.~D., Sprott D.~A. Application of
likelihood methods to models involving large numbers of parameters (with
discussion)~// Journal of the Royal Statistical Society. Series~B. 1970.
Vol.~32, №~2. P.~175--208. DOI: 10.1111/j.2517-6161.1970.tb00830.x.

\bibitem{berger1999} Berger J.~O., Liseo B., Wolpert R.~L. Integrated
likelihood methods for eliminating nuisance parameters~// Statistical Science.
1999. Vol.~14, №~1. P.~1--28. DOI: 10.1214/ss/1009211804.

\bibitem{lindsay1980} Lindsay B.~G. Nuisance parameters, mixture models, and
the efficiency of partial likelihood estimators~// Philosophical Transactions
of the Royal Society of London. Series~A. 1980. Vol.~296, №~1427. P.~639--662.
DOI: 10.1098/rsta.1980.0197.

\bibitem{garwood1936} Garwood F. Fiducial limits for the Poisson distribution~//
Biometrika. 1936. Vol.~28, №~3--4. P.~437--442. DOI: 10.1093/biomet/28.3-4.437.

\end{thebibliography}

% ── Романизированный список для международных баз цитирования ──────────────
% Все использованные источники англоязычны, поэтому транслитерация не
% требуется: список воспроизводит те же записи в том же порядке.
\begin{otherlanguage}{english}
\section*{References}
\begin{enumerate}\itemsep2pt
\item Wagstaff S.~S., Jr. Divisors of Mersenne numbers.
\emph{Mathematics of Computation}, 1983, vol.~40, no.~161, pp.~385--397.
DOI: 10.1090/S0025-5718-1983-0679454-X

\item Pomerance C. Recent developments in primality testing.
\emph{The Mathematical Intelligencer}, 1981, vol.~3, no.~3, pp.~97--105.
DOI: 10.1007/BF03022861

\item Heuristics: deriving the Wagstaff Mersenne conjecture. \emph{The PrimePages:
prime number research and records}. Available at:
\url{https://t5k.org/mersenne/heuristic.html} (accessed 27.08.2026).

\item Dubner H. Generalized repunit primes.
\emph{Mathematics of Computation}, 1993, vol.~61, no.~204, pp.~927--930.
DOI: 10.1090/S0025-5718-1993-1185243-9

\item Bourdelais P. A generalized repunit conjecture. \emph{NMBRTHRY mailing
list}, 25 June 2009. Available at:
\url{https://listserv.nodak.edu/cgi-bin/wa.exe?A2=NMBRTHRY;417ab0d6.0906}
(accessed 25.08.2026).

\item OEIS Foundation Inc. \emph{The On-Line Encyclopedia of Integer
Sequences}. Available at: \url{https://oeis.org} (accessed 27.08.2026).

\item Wagstaff S.~S., Jr. Two Mersenne prime conjectures.
\emph{Journal of Integer Sequences}, 2025, vol.~28, article 25.7.2.
Available at:
\url{https://cs.uwaterloo.ca/journals/JIS/VOL28/Wagstaff/wag14.html}
(accessed 24.09.2026).

\item Dominguez J. Arithmetic bias in the distribution of Mersenne prime exponents
and the divisor structure of $p-1$. \emph{arXiv}, 2026, arXiv:2603.08994. Available
at: \url{https://arxiv.org/abs/2603.08994} (accessed 27.08.2026).

\item Neyman J., Scott E.~L. Consistent estimates based on partially consistent
observations. \emph{Econometrica}, 1948, vol.~16, no.~1, pp.~1--32.
DOI: 10.2307/1914288

\item Andersen E.~B. Asymptotic properties of conditional maximum-likelihood
estimators. \emph{Journal of the Royal Statistical Society. Series~B}, 1970,
vol.~32, no.~2, pp.~283--301. DOI: 10.1111/j.2517-6161.1970.tb00842.x

\item Jha A. The Poisson tail conjecture for primes in short intervals.
\emph{arXiv}, 2026, arXiv:2605.23014. Available at:
\url{https://arxiv.org/abs/2605.23014} (accessed 27.08.2026).

\item Great Internet Mersenne Prime Search (GIMPS). \emph{PrimeNet}.
Prime95/MPrime software, version 30.8. Available at:
\url{https://www.mersenne.org/} (accessed 27.08.2026).

\item Epstein B. Truncated life tests in the exponential case.
\emph{The Annals of Mathematical Statistics}, 1954, vol.~25, no.~3,
pp.~555--564. DOI: 10.1214/aoms/1177728723

\item Epstein B., Sobel M. Some theorems relevant to life testing from an
exponential distribution. \emph{The Annals of Mathematical Statistics}, 1954,
vol.~25, no.~2, pp.~373--381. DOI: 10.1214/aoms/1177728793

\item Bartholomew D.~J. The sampling distribution of an estimate arising in
life testing. \emph{Technometrics}, 1963, vol.~5, no.~3, pp.~361--374.
DOI: 10.1080/00401706.1963.10490104

\item Sengupta S.~N., Sinha B.~K., Bagchi S.~B. Estimation of mean life and
reliability for exponential life distribution using time censored sample data.
\emph{Statistics \& Probability Letters}, 2009, vol.~79, no.~7, pp.~899--906.
DOI: 10.1016/j.spl.2008.11.026

\item Kivel\"{a} M., Porter M.~A. Estimating interevent time distributions
from finite observation periods in communication networks.
\emph{Physical Review E}, 2015, vol.~92, no.~5, article~052813.
DOI: 10.1103/PhysRevE.92.052813

\item Strauss D., Sadler P.~M. Classical confidence intervals and Bayesian
probability estimates for ends of local taxon ranges. \emph{Mathematical
Geology}, 1989, vol.~21, no.~4, pp.~411--427. DOI: 10.1007/BF00897326

\item Marshall C.~R. Confidence intervals on stratigraphic ranges.
\emph{Paleobiology}, 1990, vol.~16, no.~1, pp.~1--10.
DOI: 10.1017/S0094837300009672

\item Solow A.~R. Inferring extinction from sighting data. \emph{Ecology},
1993, vol.~74, no.~3, pp.~962--964. DOI: 10.2307/1940821

\item Robson D.~S., Whitlock J.~H. Estimation of a truncation point.
\emph{Biometrika}, 1964, vol.~51, no.~1--2, pp.~33--39.
DOI: 10.1093/biomet/51.1-2.33

\item Liu Z., Bhandaram U., Garg N. Quantifying spatial under-reporting
disparities in resident crowdsourcing. \emph{Nature Computational Science},
2024, vol.~4, no.~1, pp.~57--65. DOI: 10.1038/s43588-023-00572-6

\item Lancaster T. The incidental parameter problem since 1948. \emph{Journal
of Econometrics}, 2000, vol.~95, no.~2, pp.~391--413.
DOI: 10.1016/S0304-4076(99)00044-5

\item Kalbfleisch J.~D., Sprott D.~A. Application of likelihood methods to
models involving large numbers of parameters (with discussion). \emph{Journal
of the Royal Statistical Society. Series~B}, 1970, vol.~32, no.~2,
pp.~175--208. DOI: 10.1111/j.2517-6161.1970.tb00830.x

\item Berger J.~O., Liseo B., Wolpert R.~L. Integrated likelihood methods for
eliminating nuisance parameters. \emph{Statistical Science}, 1999, vol.~14,
no.~1, pp.~1--28. DOI: 10.1214/ss/1009211804

\item Lindsay B.~G. Nuisance parameters, mixture models, and the efficiency of
partial likelihood estimators. \emph{Philosophical Transactions of the Royal
Society of London. Series~A}, 1980, vol.~296, no.~1427, pp.~639--662.
DOI: 10.1098/rsta.1980.0197

\item Garwood F. Fiducial limits for the Poisson distribution.
\emph{Biometrika}, 1936, vol.~28, no.~3--4, pp.~437--442.
DOI: 10.1093/biomet/28.3-4.437
\end{enumerate}
\end{otherlanguage}

% ── Сведения об авторе ────────────────────────────────────────────────────
% п. 4.9 ГОСТ Р 7.0.7-2021: имя, отчество, фамилия полностью; организация,
% город, страна; электронный адрес БЕЗ слова «e-mail»; ORCID в форме URL.
% Элементы разделяют запятыми, точку в конце e-mail и ORCID не ставят.
\section*{Информация об авторе}
% Явный перевод строки после названия страны, а не автоматический: адрес
% набран \texttt и не переносится, поэтому в английском варианте строка
% «... Irkutsk, Russia, timofei93@yandex.ru, <ORCID>» выходила за поле.
% Разрыв стоит после запятой, разделяющей элементы по п. 4.9.
\noindent\textbf{Тимофей Сергеевич Докучаев} ---
независимый исследователь, Иркутск, Россия,\\
\texttt{timofei93@yandex.ru}, \url{https://orcid.org/0009-0006-0510-5225}

\begin{otherlanguage}{english}
\section*{Information about the author}
\noindent\textbf{Timofei Sergeevich Dokuchaev} ---
independent researcher, Irkutsk, Russia,\\
\texttt{timofei93@yandex.ru}, \url{https://orcid.org/0009-0006-0510-5225}
\end{otherlanguage}

% п. 4.20.3: сведения об отсутствии конфликта интересов приводят в конце
% статьи ПОСЛЕ «Информации об авторе», на языке статьи и затем на английском.
\medskip
\noindent Автор заявляет об отсутствии конфликта интересов.

\begin{otherlanguage}{english}
\noindent The author declares no conflict of interest.
\end{otherlanguage}

\end{document}